\section{Introduction}

A CNOT circuit on $n$ labeled wires represents an element of $\GL_n(\F_2)=\SL_n(\F_2)$, and minimizing its gate count is the word-metric problem for transvection generators.  This formulation supports asymptotically efficient synthesis, architecture-aware elimination, computable lower bounds, and exact small-width optimization \cite{SaeediMarkov2013,PatelMarkovHayes2008,JiangEtAl2020,GoubaultEtAl2021,Bataille2022,BuFanJoo2025,ChristensenEtAl2025,WebsterEtAl2025}.  Recent work has also produced stronger explicit lower-bound families and established NP-completeness of exact all-to-all CNOT distance \cite{Jorgensen2026,GongYu2026,AcuavivaEtAl2026}.  The present paper studies a different comparison: the target is lower unitriangular, but the circuit is either required to respect that ordered structure or allowed to leave it temporarily.

Let $R$ be a nonzero commutative ring with identity.  A row transvection is
\[
\tau_{t,s}(a)=I+aE_{t,s},
\qquad t\ne s,
\qquad a\in R\setminus\{0\},
\]
acting by
\[
R_t\longleftarrow R_t+aR_s.
\]
The transvection is \emph{lower directed} when $s<t$.  The lower-directed gates generate $\U_n(R)$.  We compare
\[
\ellDown(A)
 =\min\{\text{number of lower-directed transvections producing }A\}
\]
with
\[
\ellAll(A)
 =\min\{\text{number of arbitrary row transvections producing }A\},
\qquad A\in\U_n(R).
\]
For a field $\F$, the first metric is intrinsic to the unipotent radical and the second is induced from the elementary group $\Egrp_n(\F)=\SL_n(\F)$.

The equality $\ellDown(A)=\ellAll(A)$ is tempting because the target and the obvious Gaussian-elimination paths are lower unitriangular.  It is false.  The failure first occurs at six wires over $\F_2$ and $\F_3$, and it is caused by a reusable temporary carrier rather than by a better ordering of the same lower gates.

\subsection*{Main results}

The first result is entirely analytic and does not require a field.

\begin{theorem}[Universal catalytic atom over rings]\label{thm:universal-atom}
For every nonzero commutative ring $R$ with identity, let
\begin{equation}\label{eq:canonical-C}
C=
\begin{pmatrix}
1&0&0&0&0&0\\
1&1&0&0&0&0\\
0&0&1&0&0&0\\
0&1&1&1&0&0\\
1&1&1&0&1&0\\
1&1&1&0&0&1
\end{pmatrix}
\in\U_6(R).
\end{equation}
Then
\[
\ellDown(C)=7,
\qquad
\ellAll(C)=6.
\]
\end{theorem}

The six-gate ambient circuit uses row $1$ as a dirty workspace.  It first adds row $2$ into row $1$, exports the carrier $e_1+e_2$, changes the carrier to $e_0+e_1+e_2$, exports it twice, and finally removes the temporary $e_2$ component.  No extra wire is introduced.  The workspace is one of the target rows and is restored to its prescribed value.

The next result is computer assisted.  All state spaces are finite, every transition is exact, and the complete source is included in the ancillary archive.

\begin{theorem}[Exact binary threshold and census]\label{thm:binary-census}
Over $\F_2$:
\begin{enumerate}[label=\textup{(\roman*)}]
\item $\ellDown(A)=\ellAll(A)$ for every $A\in\U_n(\F_2)$ when $n\le5$.
\item Among the $2^{15}=32768$ matrices in $\U_6(\F_2)$, exactly twenty satisfy $\ellDown(A)>\ellAll(A)$.
\item Every defective target satisfies
\[
\ellDown(A)=7,
\qquad
\ellAll(A)=6.
\]
No other distance pair is defective.
\end{enumerate}
\end{theorem}

The twenty targets have a uniform combinatorial explanation.

\begin{theorem}[Minimal catalytic classification]\label{thm:motif-classification}
The twenty defective targets in $\U_6(\F_2)$ are exactly the targets generated by a simple one-catalyst word having three cross interactions, one phase interaction, and a nontrivial before/after split across the phase.  The numbers arising from the three possible left/right auxiliary splits are
\[
6,
\qquad
8,
\qquad
6,
\]
for a total of twenty.
\end{theorem}

The field-independent atom and exact ternary search also give the same finite threshold over $\F_3$.

\begin{theorem}[Exact ternary threshold]\label{thm:ternary-threshold}
For $n\le5$ and every $A\in\U_n(\F_3)$,
\[
\ellDown(A)=\ellAll(A).
\]
The matrix $C\in\U_6(\F_3)$ of \cref{eq:canonical-C} has distances $(7,6)$.  Hence six is the first defective dimension over $\F_3$.
\end{theorem}

To pass from the first threshold to arbitrarily large connected examples, we introduce a multiscale rank certificate and a family that exports successive carrier states.

\begin{theorem}[Connected snapshot family]\label{thm:snapshot-family}
Let $u\ge1$, let $1\le v\le u+1$, and choose distinct phases
\[
0\le s_1<\cdots<s_v\le u.
\]
There is a connected target
\[
A(u;s_1,\ldots,s_v)\in\U_{u+v+2}(\F)
\]
with
\[
\ellDown\bigl(A(u;s_1,\ldots,s_v)\bigr)=u+2v
\]
and
\[
\ellAll\bigl(A(u;s_1,\ldots,s_v)\bigr)\le u+v+2.
\]
Hence its certified defect is at least $v-2$.
\end{theorem}

\begin{corollary}[Connected defects in every dimension]\label{cor:connected-gap}
For every field and every $n\ge7$, there is a connected target $A_n\in\U_n(\F)$ satisfying
\[
\ellDown(A_n)-\ellAll(A_n)
\ge
\left\lfloor\frac{n-5}{2}\right\rfloor.
\]
\end{corollary}

The same construction remains effective under stringent restrictions on the ambient resource.

\begin{theorem}[Two forbidden gates and weighted penalties]\label{thm:resource-robustness}
For every field and every $n\ge7$, there is a connected target $A_n\in\U_n(\F)$ such that
\[
\ellBudget{0}(A_n)-\ellBudget{2}(A_n)
\ge
\left\lfloor\frac{n-5}{2}\right\rfloor,
\]
where $\ellBudget{q}$ permits at most $q$ upper-directed gates.  More generally, if lower gates have cost $1$ and upper gates have cost $\lambda>0$, then
\[
\ellDown(A_n)-\ellWeight{\lambda}(A_n)
\ge
\left\lfloor\frac{n-1}{2}\right\rfloor-2\lambda.
\]
Thus no fixed upper-gate budget of two, and no fixed finite upper-gate penalty, suppresses catalytic escape asymptotically.
\end{theorem}

The same paths remain extremely close to the restricted subgroup.  For a matrix $X$ over a field, let $X^+$ be its strict upper part, let $r^+(X)=\rank X^+$, and let $b^+(X)$ be the largest upper displacement $j-i$ with $X_{ij}\ne0$.  If $\ellExc{r}{b}$ denotes the minimum length among circuits satisfying $r^+(X_t)\le r$ and $b^+(X_t)\le b$ at every prefix, then the connected family can be chosen so that
\[
\ellExc{0}{0}(A_n)-\ellExc{1}{1}(A_n)
\ge
\left\lfloor\frac{n-5}{2}\right\rfloor.
\]
Thus a rank-one excursion across one adjacent flag layer can create a linear saving.

At fixed dimension, however, defects occupy thin algebraic strata.

\begin{theorem}[Generic rigidity and defect codimension]\label{thm:generic-rigidity-intro}
Let $K$ be an infinite field and put $N=\binom n2$.  There is a nonempty Zariski-open subset $\Omega_n\subseteq\U_n$ with $\Omega_n(K)\ne\varnothing$ such that
\[
\ellDown(A)=\ellAll(A)=N
\qquad(A\in\Omega_n(K)).
\]
Moreover, after base change to an algebraic closure $\overline K$, for
\[
\mathcal D_{n,\ge d}(\overline K)
=\{A\in\U_n(\overline K):\DeltaDef(A)\ge d\},
\]
one has
\[
\codim_{(\U_n)_{\overline K}}
\overline{\mathcal D_{n,\ge d}(\overline K)}\ge d.
\]
\end{theorem}

\begin{corollary}[Finite-field sparsity]\label{cor:finite-field-sparsity-intro}
Let $N=\binom n2$ and $M=n(n-1)$.  For $1\le d\le N$,
\[
\Prob_{A\in\U_n(\F_q)}[\DeltaDef(A)\ge d]
\le
\min\{1,M^{N-d}q^{-d}\}.
\]
For fixed $n$, a uniformly random target therefore has defect zero with probability $1-O_n(q^{-1})$ as $q\to\infty$.
\end{corollary}

Tensor products produce a stronger asymptotic separation.

\begin{theorem}[Tensor-amplified snapshot family]\label{thm:tensor-main}
Let $A=A(u;s_1,\ldots,s_v)$ have dimension $m=u+v+2$.  For $N=m^r$,
\[
\ellDown(A^{\otimes r})
=r m^{r-1}(u+2v),
\]
while
\[
\ellAll(A^{\otimes r})
\le r m^r.
\]
Consequently,
\[
\ellDown(A^{\otimes r})-\ellAll(A^{\otimes r})
\ge
\frac{v-2}{m\log_2m}N\log_2N.
\]
\end{theorem}

\begin{corollary}[Explicit $n\log n$ distortion]\label{cor:nlogn-gap}
For
\[
D_n(\F)
=
\sup_{A\in\U_n(\F)}
\bigl(\ellDown(A)-\ellAll(A)\bigr),
\]
one has, for every field and every $n$,
\[
D_n(\F)
\ge
\frac{8}{21\log_2 21}n\log_2n
-
\frac{82}{105}n.
\]
The coefficient is attained within the snapshot family by the twenty-one-dimensional seed
\[
A(9;0,1,\ldots,9).
\]
\end{corollary}

\subsection*{Conceptual interpretation}

The ambient circuit is shorter because one forbidden-direction gate creates a carrier that is read by several later gates before it is cleaned.  The saving is therefore a reuse effect.  The relevant object is not an ancillary wire in the usual sense.  It is a temporary excursion from a flag-preserving frame space into the full frame graph.

Over $\F_2$, the carrier states used by the canonical word form six of the seven points of $\operatorname{PG}(2,2)$.  The four projective lines supported on those six points form a Pasch configuration.  The same six vectors represent the graphic matroid $M(K_4)$ over every field; this regular-matroid structure isolates the characteristic-independent core of the atom.  At the spectral level, the distance to a target is the first nonzero coefficient of an endpoint Green kernel.  For $C$, the ambient kernel begins in degree six and the restricted kernel begins in degree seven.

The higher-dimensional families reveal three additional principles.  First, an interval-tree certificate can make the restricted metric exactly additive across scales even when the target is connected.  Second, two forbidden gates can support a linear saving, even if each forbidden gate carries any fixed finite penalty.  Third, tensor products add normalized restricted complexity while ambient complexity remains subadditive.  This converts one local carrier advantage into a superlinear global defect.  Finally, the algebraic-geometric bounds show that these amplifiable atoms are exceptional at fixed dimension: a defect of size $d$ imposes at least $d$ algebraic codimensions.  The resulting picture is one of algebraic rarity coupled with asymptotic amplification.

\subsection*{Scope}

The results concern in-place reversible linear circuits on labeled rows.  They do not by themselves imply a faster general CNOT synthesis algorithm, a lower bound for unrestricted XOR directed acyclic graphs, or a result about integer multiplication.  Their content is a precise comparison between a flag-preserving gate model and the ambient transvection model.

The paper is organized as follows.  \Cref{sec:model} fixes the circuit and metric conventions and gives the Schreier-holonomy interpretation.  \Cref{sec:atom} proves the universal ring-valued atom.  \Cref{sec:catalysis} develops the catalytic normal form and counts the twenty minimal motifs.  \Cref{sec:enumeration} gives the exact binary and ternary enumerations and proves the search procedure complete.  \Cref{sec:geometry-spectrum} gives the projective, regular-matroid, and Green-kernel formulations.  \Cref{sec:higher} develops interval-tree certificates, exact cases through dimension nine, and connected families in every larger dimension.  \Cref{sec:resources} proves bounded, weighted, and microscopic-excursion theorems.  \Cref{sec:rarity} proves generic rigidity, defect-codimension bounds, and finite-field sparsity.  \Cref{sec:amplification} proves block, recursive, and tensor amplification, introduces stable tensor capacity, and derives the $n\log n$ extremal bound.  The appendices provide the exact target list, algorithms, and reproducibility data.

\subsection*{Proof and computation status}
The universal atom, catalytic identities, interval-tree certificates, connected families, resource bounds, algebraic-rarity theorems, amplification results, and stable tensor limit are proved symbolically.  Complete search is used only for the first-threshold censuses, the equality between the twenty computed defects and the twenty analytic motifs, and the exact ambient distances of the selected seven-, eight-, and nine-dimensional targets.  Every computer-assisted statement is paired with a finite-state completeness argument and a reproducible implementation in the supplement.

\section{Circuit model and metric comparison}\label{sec:model}

\subsection{Transvections and row conventions}

Fix a nonzero commutative ring $R$ with identity and an integer $n\ge2$.  Rows and columns are indexed by
\[
0,1,\ldots,n-1.
\]
For $s\ne t$ and $a\in R\setminus\{0\}$, set
\[
\tau_{t,s}(a)=I+aE_{t,s}.
\]
Left multiplication by $\tau_{t,s}(a)$ performs
\[
R_t\longleftarrow R_t+aR_s.
\]
We call $s$ the source and $t$ the target.  A circuit word
\[
w=(\tau_{t_m,s_m}(a_m),\ldots,\tau_{t_1,s_1}(a_1))
\]
produces
\[
[w]=\tau_{t_m,s_m}(a_m)\cdots\tau_{t_1,s_1}(a_1).
\]

Let
\[
S_n^\downarrow(R)
=\{\tau_{t,s}(a):0\le s<t<n,\ a\in R\setminus\{0\}\}
\]
and
\[
S_n^*(R)
=\{\tau_{t,s}(a):0\le s,t<n,\ s\ne t,\ a\in R\setminus\{0\}\}.
\]
Both sets are symmetric because
\[
\tau_{t,s}(a)^{-1}=\tau_{t,s}(-a).
\]
Let
\[
\Egrp_n(R)=\langle S_n^*(R)\rangle.
\]
When $R=\F$ is a field and $n\ge2$, elementary elimination gives
\[
\Egrp_n(\F)=\SL_n(\F).
\]
The ambient metric below is therefore the transvection word metric in $\SL_n(\F)$, viewed inside $\GL_n(\F)$.  When $\F=\F_2$, one has $\SL_n(\F_2)=\GL_n(\F_2)$, every generator is an involution, and it is exactly a CNOT gate under the usual row-action convention.

\begin{lemma}\label{lem:lower-generates}
The set $S_n^\downarrow(R)$ generates $\U_n(R)$.  Every prefix of a lower-directed circuit belongs to $\U_n(R)$.
\end{lemma}

\begin{proof}
A lower-directed transvection is lower unitriangular, and products of lower-unitriangular matrices remain lower unitriangular.  Conversely, ordinary elimination of the strict lower entries, processed in a compatible order, expresses every element of $\U_n(R)$ as a product of lower-directed transvections.
\end{proof}

\begin{definition}[Intrinsic and ambient lengths]\label{def:lengths}
For $A\in\U_n(R)$, define
\[
\ellDown(A)=|A|_{S_n^\downarrow(R)},
\qquad
\ellAll(A)=|A|_{S_n^*(R)}.
\]
The \emph{catalytic defect} is
\[
\DeltaDef(A)=\ellDown(A)-\ellAll(A)\ge0.
\]
\end{definition}

The inequality follows because $S_n^\downarrow(R)\subseteq S_n^*(R)$.  The central question is whether the inclusion can be strict on a lower-unitriangular target.

\subsection{Frame and flag interpretation}

For the geometric interpretation in this subsection, assume $R=\F$ is a field.  An invertible matrix is an ordered basis
\[
(R_0,\ldots,R_{n-1})
\]
of $\F^n$.  A transvection replaces one basis vector by its sum with a nonzero multiple of another.  Let
\[
V_i=\Span(e_0,\ldots,e_i).
\]
A matrix lies in $\U_n(\F)$ exactly when
\[
R_i\in e_i+V_{i-1}
\]
for every $i$.  Thus lower circuits move among frames compatible with the standard complete flag
\[
V_0\subset V_1\subset\cdots\subset V_{n-1}.
\]
An upper-directed gate leaves this fiber of the flag map.  A strict defect means that the flag-compatible frame subgraph is not geodesically convex in the full transvection graph.

\subsection{Schreier holonomy}

Set
\[
G=\Egrp_n(R),
\qquad
H=\U_n(R).
\]
Consider the right Schreier graph on the left cosets $H\backslash G$: a generator $s\in S_n^*(R)$ gives an edge
\[
Hg\longrightarrow Hgs.
\]
A word representing $h\in H$ traces a based loop at the coset $H$.  Its label product is the holonomy $h$.

\begin{proposition}[Holonomy-compression interpretation]\label{prop:schreier-holonomy}
For $h\in H$, the ambient length $\ellAll(h)$ is the least length of a labeled based loop in $H\backslash G$ with holonomy $h$.  The restricted length $\ellDown(h)$ is the least length among such words whose every prefix remains in $H$.  Consequently,
\[
\DeltaDef(h)
=
\text{intrinsic flag-preserving cost}
-
\text{shortest ambient holonomy-loop cost}.
\]
\end{proposition}

\begin{proof}
A word $w=s_1\cdots s_k$ ends at the coset $Hw$.  It is a based loop exactly when $w\in H$, and its label product is then the represented element of $H$.  Minimizing over all transvection words representing $h$ gives $\ellAll(h)$.  Requiring every generator to be lower directed keeps every prefix in $H$ and gives $\ellDown(h)$.
\end{proof}

A catalytic shortcut is therefore a nonconstant quotient excursion whose accumulated holonomy is more expensive to produce inside the base fiber.

\subsection{Two elementary counting lemmas}

The proofs of the universal atom use only target counts and pivot support.

\begin{lemma}[Target count]\label{lem:target-count}
Suppose a circuit starts at the identity and finishes at $A$.
\begin{enumerate}[label=\textup{(\roman*)}]
\item If row $i$ of $A$ differs from $e_i$, then row $i$ is targeted at least once.
\item If row $i$ of $A$ equals $e_i$ and is targeted, then it is targeted at least twice.
\end{enumerate}
\end{lemma}

\begin{proof}
A row changes only when it is targeted.  Every row of an invertible matrix over $R$ is unimodular.  If $aR_s=0$ for such a row, then multiplying a unimodular relation for $R_s$ by $a$ gives $a=0$.  Hence a gate with nonzero coefficient changes its target by a nonzero vector, so one targeting cannot leave that row equal to its initial value.
\end{proof}

\begin{lemma}[One-target source condition]\label{lem:one-target-source}
Let $W\le R^n$.  Suppose row $t$ starts at $e_t$, is targeted exactly once, and finishes at
\[
e_t+w,
\qquad w\in W,
\]
where at least one coordinate of $w$ is a unit of $R$.  At the moment of its unique targeting, the gate coefficient is a unit and the source row is a unit multiple of $w$.
\end{lemma}

\begin{proof}
If the unique gate is $R_t\leftarrow R_t+aR_s$, then $aR_s=w$.  Choose a coordinate $j$ for which $w_j$ is a unit.  Since $a(R_s)_j=w_j$, the factor $a$ is a unit in the commutative ring $R$.  Hence $R_s=a^{-1}w$.
\end{proof}

The second lemma is especially restrictive when $W$ is a low-dimensional source space and output rows carry unique pivot coordinates outside $W$.

\section{The universal six-row atom}\label{sec:atom}

Let $R$ be a nonzero commutative ring with identity, and let
\[
v=e_1+e_2,
\qquad
w=e_0+e_1+e_2,
\qquad
W=\Span(e_0,e_1,e_2).
\]
The target $C$ of \cref{eq:canonical-C} has rows
\begin{align*}
R_0&=e_0,
& R_1&=e_1+e_0,
& R_2&=e_2,\\
R_3&=e_3+v,
& R_4&=e_4+w,
& R_5&=e_5+w.
\end{align*}

\begin{figure}[t]
\centering
\includegraphics[width=.56\linewidth]{canonical_matrix.pdf}
\caption{The canonical target $C$.  Black cells are nonzero.  Rows $0$ and $2$ are restored exactly, while rows $3,4,5$ require two distinct additions $v$ and $w$ from the source space $W$.}
\label{fig:canonical-matrix}
\end{figure}

\subsection{A six-gate ambient word}

\begin{proposition}\label{prop:ambient-upper}
For every nonzero commutative ring $R$ with identity,
\[
\ellAll(C)\le6.
\]
\end{proposition}

\begin{proof}
Apply the following six gates:
\begin{align}
R_1&\leftarrow R_1+R_2,\label{gate:a1}\\
R_3&\leftarrow R_3+R_1,\label{gate:a2}\\
R_1&\leftarrow R_1+R_0,\label{gate:a3}\\
R_4&\leftarrow R_4+R_1,\label{gate:a4}\\
R_5&\leftarrow R_5+R_1,\label{gate:a5}\\
R_1&\leftarrow R_1-R_2.\label{gate:a6}
\end{align}
After \cref{gate:a1}, row $1$ contains $v$.  Gate \cref{gate:a2} copies $v$ to row $3$.  Gate \cref{gate:a3} changes the carrier to $w$, which is copied to rows $4$ and $5$ by \cref{gate:a4,gate:a5}.  Gate \cref{gate:a6} removes the temporary $e_2$ component, leaving $e_1+e_0$.  Every row now agrees with $C$.
\end{proof}

\begin{figure}[t]
\centering
\includegraphics[width=.98\linewidth]{carrier_timeline.pdf}
\caption{Carrier timeline for the six-gate ambient word.  The two upper-directed operations are the opening and closing gates on row $1$.  The state $e_1+e_2$ is exported once and $e_0+e_1+e_2$ is exported twice.}
\label{fig:carrier-timeline}
\end{figure}

\subsection{A seven-gate lower word}

\begin{proposition}\label{prop:lower-upper}
For every nonzero commutative ring $R$ with identity,
\[
\ellDown(C)\le7.
\]
\end{proposition}

\begin{proof}
The lower-directed word
\begin{align*}
R_1&\leftarrow R_1+R_0,\\
R_2&\leftarrow R_2+R_1,\\
R_3&\leftarrow R_3-R_0,\\
R_3&\leftarrow R_3+R_2,\\
R_4&\leftarrow R_4+R_2,\\
R_5&\leftarrow R_5+R_2,\\
R_2&\leftarrow R_2-R_1
\end{align*}
produces $C$.  After the first two gates, row $2$ contains $w$.  Rows $4$ and $5$ copy that state.  The two gates on row $3$ cancel the unwanted $e_0$ and leave $v$.  The final gate restores row $2$.
\end{proof}

\subsection{Restricted lower bound}

\begin{proposition}\label{prop:lower-lower-bound}
For every nonzero commutative ring $R$ with identity,
\[
\ellDown(C)\ge7.
\]
\end{proposition}

\begin{proof}
Rows $1,3,4,5$ differ from their initial values, so they must each be targeted.  Assume for contradiction that a lower-directed word of length at most six produces $C$.

First suppose row $2$ is never targeted.  Row $0$ remains $e_0$, row $2$ remains $e_2$, and every state ever assumed by row $1$ lies in $\Span(e_0,e_1)$ because the only possible source below row $1$ is row $0$.  Every output row retains coefficient $1$ on its own pivot coordinate in a lower-directed circuit, so no output row can lie in $W$.  Hence no available source row is a unit multiple of either
\[
v=e_1+e_2
\quad\text{or}\quad
w=e_0+e_1+e_2.
\]
By \cref{lem:one-target-source}, each of rows $3,4,5$ must be targeted at least twice.  Together with the required targeting of row $1$, this gives at least seven gates.

Now suppose row $2$ is targeted.  Since its final value is $e_2$, \cref{lem:target-count} forces at least two targetings of row $2$.  A six-gate word must therefore target row $2$ exactly twice and each of rows $1,3,4,5$ exactly once.  Each output row $3,4,5$ must receive its complete addition in its unique targeting.  Row $0$ cannot supply $v$ or $w$, and row $1$ never contains an $e_2$ component.  Therefore row $2$ must supply all three outputs, up to unit scaling.  But two targetings of row $2$ expose at most one noninitial intermediate state before row $2$ returns to $e_2$, while the outputs require both distinct states $v$ and $w$.  This is impossible.
\end{proof}

\subsection{Ambient lower bound}

\begin{proposition}\label{prop:ambient-lower-bound}
For every nonzero commutative ring $R$ with identity,
\[
\ellAll(C)\ge6.
\]
\end{proposition}

\begin{proof}
Assume that an ambient word of length at most five produces $C$.  Rows $1,3,4,5$ must each be targeted.  Rows $0$ and $2$ are unchanged.  If either were targeted, it would require at least two targetings, giving at least six gates in total.  Thus rows $0$ and $2$ are never targeted.  If all four changed rows are targeted once, each has only its initial and final state.  No such source is a unit multiple of $v=e_1+e_2$ or $w=e_0+e_1+e_2$, so the output rows cannot be completed.  In the only remaining case, exactly one of rows $1,3,4,5$ is targeted twice and the other three are targeted once.

Project the rows to the quotient module $R^6/W$.  Rows $0,1,2$ have zero image, while rows $3,4,5$ begin and end at the three distinct quotient basis vectors of the free quotient module.  An output row targeted exactly once must therefore use a source lying in $W$, by \cref{lem:one-target-source}.

Suppose first that row $1$ is the twice-targeted row.  The three output rows are each targeted once.  Rows $0$ and $2$ supply only scalar multiples of $e_0$ and $e_2$.  Row $1$ begins at $e_1$, has one intermediate state, and ends at $e_1+e_0$.  Thus at most one of the two required unit-multiple classes of $v$ and $w$ can appear as an available source.  The three outputs require both, a contradiction.

Suppose instead that an output row, say row $r\in\{3,4,5\}$, is targeted twice.  The other two outputs and row $1$ are targeted once.  We claim that row $r$ cannot lie in $W$ between its two targetings.  To cancel its unique quotient pivot $e_r$ at the first targeting, the source would already have to contain an $e_r$ component.  Before row $r$ is first targeted, the only way another row can acquire such a component is to be targeted using row $r$ as source.  That other row is targeted only once and would then retain an unwanted $e_r$ component in its final value.  Therefore row $r$ never becomes a $W$-valued carrier.  The only $W$-valued source states are rows $0$ and $2$ and the two states of singly targeted row $1$, namely its initial and final states.  None supplies both $v$ and $w$.  The two singly targeted outputs cannot be completed, again a contradiction.
\end{proof}

\begin{proof}[Proof of \cref{thm:universal-atom}]
Combine \cref{prop:ambient-upper,prop:lower-upper,prop:lower-lower-bound,prop:ambient-lower-bound}.
\end{proof}

\begin{remark}
The proof uses only unimodular rows, unit coordinates, and pivot support.  It does not use division or a characteristic assumption.  In characteristic two the closing subtraction equals addition.
\end{remark}

\begin{corollary}[Failure of geodesic convexity]\label{cor:nonconvex}
For every field $\F$ and every $n\ge6$, the lower-unitriangular frame subgraph is not geodesically convex in the transvection Cayley graph of $\Egrp_n(\F)=\SL_n(\F)$.
\end{corollary}

\begin{proof}
Embed $C$ as $C\oplus I_{n-6}$.  A six-gate ambient word acts on the first six rows.  In a lower-directed circuit, rows with index at least six cannot be used as sources for targets among the first six rows, so the restricted lower bound of seven persists.
\end{proof}

For the remainder of the paper, unless explicitly stated otherwise, the coefficient domain is a field $\F$.

\section{Catalytic normal forms and the six-role threshold}\label{sec:catalysis}

The universal atom is one instance of a root-conjugation mechanism.  This section isolates that mechanism and explains the count twenty.

\subsection{Root conjugation identities}

Let $p<q<y$ and $x<p<q$.  Direct multiplication gives
\begin{equation}\label{eq:read-conjugation}
\tau_{p,q}(-a)\,\tau_{y,p}(b)\,\tau_{p,q}(a)
=
\tau_{y,p}(b)\,\tau_{y,q}(ab),
\end{equation}
and
\begin{equation}\label{eq:update-conjugation}
\tau_{p,q}(-a)\,\tau_{q,x}(b)\,\tau_{p,q}(a)
=
\tau_{q,x}(b)\,\tau_{p,x}(-ab).
\end{equation}
A single forbidden upper transvection can therefore turn one lower interaction into two lower effects.  Interactions of the forms $x\to p$ and $q\to y$ commute through the excursion without duplication.  We call the duplicated interactions \emph{cross interactions} and the unduplicated interactions \emph{phase interactions}.

\begin{definition}[Simple one-catalyst word]\label{def:simple-catalyst}
A simple one-catalyst word has the form
\[
\tau_{p,q}(-a)\,h_m\cdots h_1\,\tau_{p,q}(a),
\qquad p<q,
\]
where every $h_i$ is lower directed and is one of
\[
\tau_{q,x}(b),\quad \tau_{p,x}(b),\quad
\tau_{y,p}(b),\quad \tau_{y,q}(b)
\]
with $x<p<q<y$.  Auxiliary rows are distinct unless stated otherwise.
\end{definition}

The circuit is read from right to left in the matrix product and from left to right as a gate list.  The opening and closing operations temporarily place the $q$ component in row $p$.

\begin{proposition}[Cross-interaction accounting]\label{prop:cross-accounting}
Let a simple one-catalyst word have $m$ middle interactions, of which $c$ are cross interactions.  Its target admits a lower-directed realization of length at most
\[
m+c.
\]
Consequently, if the ambient word of length $m+2$ is strictly shorter than every lower word, then
\[
c\ge3.
\]
\end{proposition}

\begin{proof}
Conjugate each middle generator using \cref{eq:read-conjugation,eq:update-conjugation}.  A phase interaction remains one lower transvection, while a cross interaction becomes a product of two lower transvections.  This gives length at most $m+c$.  If $c\le2$, then $m+c\le m+2$, so the ambient excursion cannot be strictly better.
\end{proof}

Three consumers are therefore the first possible point at which the cost two of opening and closing can be amortized.

\subsection{Why an all-cross word is insufficient}

Within the distinct-leaf normal form, cross count alone is not enough.

\begin{lemma}[All-cross mirror]\label{lem:all-cross-mirror}
Suppose a simple one-catalyst word contains only cross interactions, with left inputs $X$ and right outputs $Y$.  Then its target has a lower-directed realization with the same number of gates.
\end{lemma}

\begin{proof}
Write the opening coefficient as $a\in\F^\times$, the left-update coefficients as $b_x$, and the right-read coefficients as $c_y$.  Since the all-cross middle gates commute past one another, the ambient gate list may be written
\[
R_p\leftarrow R_p+aR_q;
\quad
[R_q\leftarrow R_q+b_xR_x]_{x\in X};
\quad
[R_y\leftarrow R_y+c_yR_p]_{y\in Y};
\quad
R_p\leftarrow R_p-aR_q.
\]
Put $B=\sum_{x\in X}b_xe_x$.  Its final nontrivial rows are
\[
R_p=e_p-aB,
\qquad
R_q=e_q+B,
\qquad
R_y=e_y+c_ye_p+ac_ye_q
\quad(y\in Y).
\]
The following lower-directed gate list produces the same rows:
\[
R_q\leftarrow R_q+a^{-1}R_p;
\quad
[R_y\leftarrow R_y+ac_yR_q]_{y\in Y};
\quad
[R_p\leftarrow R_p-ab_xR_x]_{x\in X};
\quad
R_q\leftarrow R_q-a^{-1}R_p.
\]
Indeed, after the first gate $R_q=e_q+a^{-1}e_p$, so the output gates create exactly $c_ye_p+ac_ye_q$.  The source updates then leave $R_p=e_p-aB$, and the final gate changes $R_q$ to $e_q+B$.  Both words have $|X|+|Y|+2$ gates.
\end{proof}

Thus a minimal strict shortcut needs both positive catalytic excess and a phase switch that makes temporal order matter.

\begin{proposition}[Six-role threshold in the simple model]\label{prop:six-role}
A strict shortcut in the distinct-leaf simple one-catalyst model requires at least six row roles:
\[
2\text{ catalyst rows}
+
3\text{ cross interactions}
+
1\text{ phase interaction}.
\]
The bound is attained by the matrix $C$.
\end{proposition}

\begin{proof}
By \cref{prop:cross-accounting}, at least three cross interactions are required.  By \cref{lem:all-cross-mirror}, at least one phase interaction is required.  With distinct leaves these use four auxiliary rows in addition to the two catalyst rows.  The word of \cref{prop:ambient-upper} has exactly this form.
\end{proof}

\begin{figure}[t]
\centering
\includegraphics[width=.82\linewidth]{catalytic_motif.pdf}
\caption{The minimal one-catalyst motif.  Cross interactions reuse the opened carrier.  The phase gate separates consumers that observe two different carrier states and prevents the all-cross mirror.}
\label{fig:catalytic-motif}
\end{figure}

\subsection{Enumeration of the minimal motifs}

Fix six ordered rows.  The catalyst pair must be adjacent in a minimal six-role word.  Write
\[
p<q=p+1,
\]
let $L$ be the set of rows left of $p$, and let $R$ be the set of rows right of $q$.  Then
\[
|L|+|R|=4,
\qquad
|L|,|R|>0.
\]
There are two orientations.

\paragraph{Left phase.}
Choose a phase leaf $x\in L$.  The remaining $|L|-1$ left leaves interact by $u\to q$.  Every right leaf is a cross read $p\to y$.  To make the phase effective, the right leaves are divided into a nonempty set read before $x\to p$ and a nonempty set read after it.  The number of targets is
\[
|L|\bigl(2^{|R|}-2\bigr).
\]

\paragraph{Right phase.}
Choose a phase leaf $y\in R$.  Every other right leaf is a cross read.  The left updates into $q$ are divided nontrivially before and after the phase gate $q\to y$.  The number is
\[
|R|\bigl(2^{|L|}-2\bigr).
\]

\begin{proposition}[Motif count]\label{prop:motif-count}
For a split $(L,R)=(|L|,|R|)$ with $L+R=4$, the number of minimal phase-separated targets is
\[
N_{L,R}
=L(2^R-2)+R(2^L-2).
\]
For the three possible splits,
\[
N_{1,3}=6,
\qquad
N_{2,2}=8,
\qquad
N_{3,1}=6.
\]
Hence there are twenty targets.
\end{proposition}

\begin{proof}
The preceding left-phase and right-phase constructions are disjoint because the phase row lies on opposite sides of the catalyst.  Within each orientation, the final support records the phase leaf and the before/after subset, so distinct choices give distinct targets.  Summing gives the formula and the stated values.
\end{proof}

\begin{figure}[t]
\centering
\includegraphics[width=.70\linewidth]{motif_count.pdf}
\caption{The analytic count of the twenty minimal catalytic targets.}
\label{fig:motif-count}
\end{figure}

Every target counted by \cref{prop:motif-count} has an explicit six-gate ambient word and a seven-gate lower word obtained by eliminating the three cross interactions.  The nontrivial statement is that there are no other defects in dimension six and that none of these targets has a lower word of length six.  That is the exact classification established in \cref{sec:enumeration}.

\section{Exact finite enumeration}\label{sec:enumeration}

This section describes the computer-assisted part of the paper.  The computations use integer encodings only.  No floating-point decisions occur.

\subsection{Restricted breadth-first search}

Over $\F_2$, an element of $\U_n(\F_2)$ is determined by its $n(n-1)/2$ strict lower entries.  For $n=6$ there are $2^{15}=32768$ states.  Breadth-first search from the identity using the fifteen lower transvections gives every restricted distance and the restricted diameter.

For $n=6$, the restricted diameter is eight.  Therefore every lower-unitriangular target has ambient distance at most eight.  This makes a radius-four meet-in-the-middle search complete.

\subsection{Ambient radius-four ball and right-coset pairing}

Let
\[
B_4=\{g\in\GL_6(\F_2):\ellAll(g)\le4\}.
\]
Exact breadth-first search gives
\[
|B_4|=110626.
\]
Group $B_4$ by right cosets of $\U_6(\F_2)$.  If $a,b$ lie in the same right coset, then
\[
b^{-1}a\in\U_6(\F_2),
\]
and
\[
\ellAll(b^{-1}a)\le\ellAll(b)+\ellAll(a)\le8.
\]
The implementation uses a canonical right-coset key obtained by reducing matrix columns from right to left.  Right multiplication by a lower-unitriangular matrix changes column $j$ only by a combination of columns with larger index, so this reduction is constant on right cosets and separates them.

\begin{lemma}[Meet-in-the-middle completeness]\label{lem:mitm-complete}
For every $u\in\U_6(\F_2)$,
\[
\ellAll(u)
=
\min\{\ellAll(a)+\ellAll(b):a,b\in B_4,\ b^{-1}a=u\}.
\]
\end{lemma}

\begin{proof}
Every pair on the right gives a word for $u$, so the right-hand side is at least $\ellAll(u)$.  Conversely, $\ellAll(u)\le8$.  Take a shortest path from the identity to $u$ and split it at a midpoint.  The prefix endpoint $a$ has distance at most four.  The inverse of the remaining suffix is an element $b$ of distance at most four, and $u=b^{-1}a$.  Hence the minimum is at most $\ellAll(u)$.
\end{proof}

The radius-four ball splits into $6948$ right-coset buckets.  The largest bucket has $3775$ elements.  The computation checks $45178326$ ordered pairs and recovers an ambient distance for every one of the $32768$ targets.

\subsection{The exact binary census}

\begin{table}[t]
\centering
\caption{Complete distance distribution in $\U_6(\F_2)$.}
\label{tab:census}
\begin{tabular}{@{}rrr@{}}
\toprule
$\ellDown$ & $\ellAll$ & number of targets\\
\midrule
0&0&1\\
1&1&15\\
2&2&125\\
3&3&720\\
4&4&2914\\
5&5&7851\\
6&6&12656\\
7&6&20\\
7&7&7926\\
8&8&540\\
\bottomrule
\end{tabular}
\end{table}

\begin{figure}[t]
\centering
\includegraphics[width=.72\linewidth]{distance_census.pdf}
\caption{Complete binary dimension-six census.  The only off-diagonal cell contains the twenty targets with distance pair $(7,6)$.}
\label{fig:distance-census}
\end{figure}

\begin{proof}[Proof of \cref{thm:binary-census}]
For $n=2,3,4,5$, complete restricted breadth-first search and complete ambient meet-in-the-middle search return zero defects.  The state counts are respectively $2,8,64,1024$.  For $n=6$, \cref{lem:mitm-complete} proves that the ambient calculation is exhaustive.  The resulting distribution is \cref{tab:census}; its counts sum to $32768$, no target is missing, and the only strict pair is $(7,6)$ with multiplicity twenty.
\end{proof}

\begin{proof}[Proof of \cref{thm:motif-classification}]
The analytic generator of \cref{prop:motif-count} produces twenty distinct targets without reference to the distance census.  Exact set comparison shows that these twenty targets are precisely the twenty off-diagonal targets in \cref{tab:census}.  Each therefore has distances $(7,6)$, and no other dimension-six target is defective.
\end{proof}

\subsection{Ternary and additional field checks}

For $\F_3$, every nonzero scaled transvection is included.  Complete searches give the data in \cref{tab:small-fields}.

\begin{table}[t]
\centering
\caption{Complete no-defect computations below dimension six.}
\label{tab:small-fields}
\begin{tabular}{@{}crrr@{}}
\toprule
field & $n$ & $|\U_n(\F)|$ & restricted diameter\\
\midrule
$\F_2$&2&2&1\\
$\F_2$&3&8&2\\
$\F_2$&4&64&4\\
$\F_2$&5&1024&6\\
\addlinespace
$\F_3$&2&3&1\\
$\F_3$&3&27&3\\
$\F_3$&4&729&5\\
$\F_3$&5&59049&7\\
\bottomrule
\end{tabular}
\end{table}

The $n=5$ ternary computation uses radius-three and radius-four balls, checks $38415205$ right-coset pairs, and recovers all $59049$ targets.  It finds no defect.  The field-independent \cref{thm:universal-atom} supplies a defect at $n=6$, proving \cref{thm:ternary-threshold}.  Separate exact searches over $\F_3$ and $\F_5$ verify the distances of the canonical atom directly.  The $\F_5$ calculation is an audit of the atom, not a complete threshold computation below dimension six.

\subsection{Verification status}

The analytic proofs of \cref{thm:universal-atom,cor:linear-gap} do not depend on computation.  The claims in \cref{thm:binary-census,thm:motif-classification,thm:ternary-threshold} depend on exact finite enumeration.  The source archive contains:
\begin{itemize}
\item the C++ programs for every breadth-first and meet-in-the-middle search;
\item the Python program that independently generates the twenty motifs and compares sets;
\item symbolic checks for the bounded-resource formulas and the totally unimodular $M(K_4)$ representation;
\item exact stdout files and JSON packet data;
\item SHA-256 checksums;
\item build and reproduction commands.
\end{itemize}
The programs were rebuilt from source and rerun independently when this manuscript was assembled; the outputs matched byte for byte.

\section{Projective, matroidal, and spectral structure}\label{sec:geometry-spectrum}

\subsection{The flagged Pasch packet}

Over $\F_2$, write
\[
x=e_0,
\qquad
p=e_1,
\qquad
q=e_2.
\]
The canonical ambient word uses the six nonzero vectors
\[
x,
\quad p,
\quad q,
\quad p+q,
\quad p+q+x,
\quad p+x
\]
inside $\Span(x,p,q)\cong\F_2^3$.  The only missing nonzero vector is $x+q$.

In $\operatorname{PG}(2,2)$, three distinct points form a line exactly when their vector representatives sum to zero.  The six points support the four lines
\[
\{x,p,x+p\},
\quad
\{x,p+q,x+p+q\},
\quad
\{p,q,p+q\},
\quad
\{x+p,q,x+p+q\}.
\]
Every point lies on exactly two of these lines.  This is the Pasch configuration, equivalently the Fano plane with one point and its three incident lines removed from the full incidence structure.

\begin{figure}[t]
\centering
\includegraphics[width=.65\linewidth]{pasch_configuration.pdf}
\caption{A real drawing of the four-line Pasch incidence structure.  The labels are the binary point encodings $1,2,3,4,6,7$; point $5$ is the missing Fano point.}
\label{fig:pasch}
\end{figure}

The configuration is flagged because the rows have fixed roles.  Points $x$ and $q$ are stationary sources, $p$ is the carrier, and the remaining points occur as carrier states.  Forgetting those roles loses the circuit ordering information.

\subsection{The regular matroid \texorpdfstring{$M(K_4)$}{M(K4)}}

The six carrier states admit a field-independent representation.  In the ordered basis $(x,p,q)$, place them as the columns of
\begin{equation}\label{eq:k4-matrix}
B=
\begin{pmatrix}
1&0&0&1&0&1\\
0&1&0&1&1&1\\
0&0&1&0&1&1
\end{pmatrix},
\end{equation}
corresponding respectively to
\[
x,
\quad p,
\quad q,
\quad x+p,
\quad p+q,
\quad x+p+q.
\]

\begin{proposition}[Regular-matroid core]\label{prop:regular-matroid}
Every square minor of $B$ belongs to $\{0,\pm1\}$, so $B$ is totally unimodular.  Hence the column-dependence matroid is the same over every field.  Under the edge labeling
\[
\begin{array}{c|cccccc}
\text{state}&p&x&x+p&p+q&q&x+p+q\\
\hline
K_4\text{ edge}&12&13&23&14&24&34,
\end{array}
\]
it is the graphic matroid $M(K_4)$.
\end{proposition}

\begin{proof}
A direct calculation of all square minors of $B$ gives only $0$, $1$, and $-1$.  Thus every subdeterminant has characteristic-independent vanishing, and a column set is independent over one field exactly when it is independent over every field.  The four minimal dependent triples are
\[
\{x,p,x+p\},
\quad
\{p,q,p+q\},
\quad
\{x,p+q,x+p+q\},
\quad
\{q,x+p,x+p+q\}.
\]
Under the displayed edge labeling, these are precisely the four triangular cycles of $K_4$.  They determine the rank-three cycle matroid $M(K_4)$, which is regular because it has a totally unimodular representation; see, for example, \cite{Oxley2011}.
\end{proof}

\begin{figure}[t]
\centering
\includegraphics[width=.72\linewidth]{regular_matroid_k4.pdf}
\caption{The six carrier states as the edges of $K_4$.  Each triangular cycle is one of the four three-element dependencies in the Pasch configuration.  The resulting regular matroid is representable over every field.}
\label{fig:k4-matroid}
\end{figure}

The field-independent circuit theorem is stronger than matroid representability alone, because it also depends on the ordering and the designated carrier role.  Nevertheless, \cref{prop:regular-matroid} explains why the local dependence pattern survives unchanged in every characteristic.

\subsection{Distance as first arrival}

Let $G$ be a finite group and $S=S^{-1}$ a symmetric generating set.  On $\ell^2(G)$, let $\lambda$ be the left regular representation,
\[
\lambda(s)\delta_h=\delta_{sh},
\]
and define
\[
A_S=\sum_{s\in S}\lambda(s).
\]
For $g\in G$, define the endpoint Green series
\[
\mathcal G_{S,g}(z)
=
\left\langle
\delta_g,
(I-zA_S)^{-1}\delta_e
\right\rangle.
\]

\begin{proposition}[First-arrival identity]\label{prop:first-arrival}
The coefficient of $z^k$ in $\mathcal G_{S,g}(z)$ is the number of length-$k$ words in $S$ with product $g$.  Hence
\[
d_S(e,g)=\ord_{z=0}\mathcal G_{S,g}(z).
\]
\end{proposition}

\begin{proof}
Formally,
\[
(I-zA_S)^{-1}=\sum_{k\ge0}z^kA_S^k.
\]
The matrix coefficient $\langle\delta_g,A_S^k\delta_e\rangle$ counts $k$-step walks from $e$ to $g$, which are exactly length-$k$ words with product $g$.  The first nonzero coefficient occurs at the word distance.
\end{proof}

If
\[
A_S=\sum_\lambda \lambda\Pi_\lambda
\]
is the spectral resolution, then
\[
\mathcal G_{S,g}(z)
=
\sum_\lambda
\frac{\langle\delta_g,\Pi_\lambda\delta_e\rangle}{1-\lambda z}.
\]
Thus eigenvalues alone do not determine first arrival.  The endpoint matrix coefficients of the spectral projectors are essential.

For the canonical target $C$, exact word counts begin
\begin{align*}
\mathcal G_{\downarrow,C}(z)
&=
543z^7+25080z^8+510703z^9+10824940z^{10}+\cdots,\\
\mathcal G_{*,C}(z)
&=
2z^6+776z^7+39454z^8+\cdots.
\end{align*}
The metric defect is the valuation drop
\[
\ord_0\mathcal G_{\downarrow,C}
-
\ord_0\mathcal G_{*,C}
=1.
\]

\begin{figure}[t]
\centering
\includegraphics[width=.78\linewidth]{green_coefficients.pdf}
\caption{Exact endpoint word counts for the canonical target.  The ambient walk first arrives at length six.  The restricted walk first arrives at length seven.}
\label{fig:green-coefficients}
\end{figure}

Across the twenty defects, the truncated endpoint data split into five Green-signature classes, each of multiplicity four.  The signatures are listed in \cref{tab:green-signatures}.

\begin{table}[t]
\centering
\caption{The five endpoint Green-signature packets.  Each row occurs for four defective targets.}
\label{tab:green-signatures}
\resizebox{\textwidth}{!}{%
\begin{tabular}{@{}rrrrrrrr@{}}
\toprule
&\multicolumn{4}{c}{restricted coefficients}&\multicolumn{3}{c}{ambient coefficients}\\
packet&$[z^7]$&$[z^8]$&$[z^9]$&$[z^{10}]$&$[z^6]$&$[z^7]$&$[z^8]$\\
\midrule
1&326&15430&316118&7576819&4&570&23318\\
2&326&15536&321870&7727183&4&570&23318\\
3&543&25080&510703&10824940&2&776&39454\\
4&543&25639&511524&10984988&2&776&39454\\
5&543&26091&511819&11121486&2&776&39454\\
\bottomrule
\end{tabular}}
\end{table}

\section{Interval-tree certificates and connected higher-dimensional families}\label{sec:higher}

From this point through the rank and tensor arguments, the coefficient domain is a field $\F$.  The six-row atom is the first defect, but it is not an isolated low-dimensional accident.  This section gives a field-independent family of connected targets in every dimension at least seven.  The lower bound is certified by a multiscale rank potential adapted to the ordered flag.

\subsection{Interval-tree rank potential}

\begin{definition}[Interval tree]
An \emph{interval tree} on the ordered coordinates $0,1,\ldots,n-1$ is a rooted binary tree whose leaves occur in that order and whose every node is a contiguous interval.  If an internal node $v$ splits its interval into consecutive left and right children $L_v$ and $R_v$, define
\begin{equation}\label{eq:tree-potential}
\Phi_T(A)
=
\sum_{v\in\operatorname{Int}(T)}
\rank A[R_v,L_v],
\qquad A\in\U_n(\F).
\end{equation}
\end{definition}

\begin{lemma}[One-cut locality]\label{lem:one-cut-locality}
Let $X\in\U_n(\F)$ and let $X'$ be obtained from $X$ by one lower-directed transvection.  Then
\[
\bigl|\Phi_T(X')-\Phi_T(X)\bigr|\le1.
\]
\end{lemma}

\begin{proof}
Suppose the gate is $s\to t$ with $s<t$.  Let $v_0$ be the lowest node of $T$ whose two children separate $s$ and $t$.  The order of the leaves forces $s\in L_{v_0}$ and $t\in R_{v_0}$.  At this one node, a single row of the cross block $X[R_{v_0},L_{v_0}]$ changes, so its rank changes by at most one.

At every ancestor of $v_0$, the two rows lie in the same child.  If that child is the right child, the gate is an internal row operation on the relevant cross block and preserves rank; if it is the left child, the cross block is unchanged.  At a descendant containing $t$ but not $s$, every column of the descendant interval lies strictly to the right of $s$.  Since $X$ remains lower triangular, row $s$ vanishes on those columns, so the descendant block is unchanged.  All other descendants are also unchanged.  Hence only the term at $v_0$ can vary.
\end{proof}

\begin{corollary}[Interval-tree lower bound]\label{cor:tree-lower-bound}
For every $A\in\U_n(\F)$ and every interval tree $T$,
\[
\ellDown(A)\ge\Phi_T(A).
\]
\end{corollary}

\begin{proof}
The identity has potential zero.  Apply \cref{lem:one-cut-locality} along any lower-directed word from the identity to $A$.
\end{proof}

\begin{definition}[Tree-tight target]\label{def:tree-tight}
A lower-unitriangular target $A$ is \emph{tree tight} if some interval tree $T$ satisfies
\[
\Phi_T(A)=\ellDown(A).
\]
\end{definition}

Tree tightness is stronger than possessing a rank lower bound.  It says that one ordered hierarchy accounts for the entire restricted word length.  The family below is tree tight by construction.

\subsection{The snapshot family}

Fix integers
\[
u\ge1,
\qquad
1\le v\le u+1,
\]
and choose distinct phases
\[
0\le s_1<s_2<\cdots<s_v\le u.
\]
Order the coordinates as
\[
x_1,\ldots,x_u,\ p,\ q,\ t_1,\ldots,t_v.
\]
For $0\le s\le u$, put
\begin{equation}\label{eq:prefix-state}
w_s=e_p+\sum_{i=1}^{s}e_{x_i}.
\end{equation}
Define
\begin{equation}\label{eq:snapshot-target}
A(u;s_1,\ldots,s_v)\in\U_{u+v+2}(\F)
\end{equation}
by
\[
R_{x_i}=e_{x_i},
\qquad
R_p=w_u,
\qquad
R_q=e_q,
\]
and
\begin{equation}\label{eq:snapshot-output}
R_{t_j}=e_{t_j}+e_q+w_{s_j}
\qquad(1\le j\le v).
\end{equation}
The output rows record selected intermediate versions of the carrier $p$.

\begin{proof}[Proof of \cref{thm:snapshot-family}]
For the ambient construction, first apply
\[
q\to p.
\]
Sweep through the updates
\[
x_1\to p,\ldots,x_u\to p.
\]
Whenever phase $s_j$ is reached, apply
\[
p\to t_j.
\]
Finally apply $q\to p$ with coefficient $-1$.  At phase $s$, the carrier is $e_q+w_s$, so every output has the required value, and the final gate removes $e_q$ from row $p$.  The word has length
\[
1+u+v+1=u+v+2.
\]
Thus
\[
\ellAll\bigl(A(u;s_1,\ldots,s_v)\bigr)\le u+v+2.
\]

For a lower word, perform the same $u$ updates without opening the catalyst.  At each selected phase, use the two gates
\[
p\to t_j,
\qquad
q\to t_j.
\]
This gives length $u+2v$.

It remains to prove optimality in the lower model.  At the root, split
\[
\{x_1,\ldots,x_u,p\}
\quad\big|\quad
\{q,t_1,\ldots,t_v\}.
\]
The nonzero rows of the root cross block are $w_{s_1},\ldots,w_{s_v}$.  Distinct prefix vectors are linearly independent: in a linear relation, inspect successively a coordinate occurring in $w_{s_j}$ but in none of the earlier prefixes, starting from the largest phase.  Hence the root contributes $v$.

Inside the left child, use the comb that splits off $x_1$, then $x_2$, and so on while retaining $p$ in the remaining interval.  Every one of the $u$ cuts has rank one because the final row $p$ contains every $x_i$.  Inside the right child, repeatedly split off the last output while retaining $q$ in the left interval.  Each of the $v$ cuts has rank one because every output contains $q$.  Therefore
\[
\Phi_T\bigl(A(u;s_1,\ldots,s_v)\bigr)
=v+u+v=u+2v.
\]
The interval-tree lower bound matches the displayed lower circuit.
\end{proof}

\begin{figure}[t]
\centering
\includegraphics[width=.90\linewidth]{snapshot_family.pdf}
\caption{The connected snapshot family.  One temporary upper dependency lets the ambient circuit export $v$ selected prefix states.  A lower-only circuit must insert the $q$ component separately into every output.}
\label{fig:snapshot-family}
\end{figure}

\begin{remark}[Support connectivity]
For $u,v\ge1$, the undirected support graph of $A(u;s_1,\ldots,s_v)-I$ is connected.  Row $p$ joins all update coordinates $x_i$, while every output joins both $p$ and $q$.  The separation below therefore does not arise from a direct sum of independent six-row gadgets.
\end{remark}

\subsection{Exact higher binary examples}

The interval tree gives the exact restricted distance in all dimensions.  For the first three higher targets, a separate bidirectional search also proves that the displayed ambient construction is optimal.

\begin{theorem}[Exact binary cases in dimensions seven through nine]\label{thm:higher-exact}
Over $\F_2$, let
\[
H_7=A(2;0,1,2),
\qquad
H_8=A(3;0,1,2),
\qquad
H_9=A(3;0,1,2,3).
\]
Then
\[
(\ellDown(H_7),\ellAll(H_7))=(8,7),
\]
\[
(\ellDown(H_8),\ellAll(H_8))=(9,8),
\]
and
\[
(\ellDown(H_9),\ellAll(H_9))=(11,9).
\]
\end{theorem}

\begin{proof}
The restricted values follow from \cref{thm:snapshot-family}.  The ambient words have lengths $7$, $8$, and $9$.  To exclude shorter words, the supplement computes exact ambient balls and checks that the two relevant balls do not intersect.  The certificate sizes are
\begin{center}
\begin{tabular}{@{}ccll@{}}
\toprule
$n$&tested split&left ball&right ball\\
\midrule
$7$&$3+3$&$23276$&$23276$\\
$8$&$3+4$&$52277$&$1174097$\\
$9$&$4+4$&$2994949$&$2994949$\\
\bottomrule
\end{tabular}
\end{center}
Every intersection is empty.  Since the generating sets are symmetric, an empty radius-$r$ versus radius-$s$ intersection rules out every word of length at most $r+s$.
\end{proof}

These are exact higher-dimensional examples, not merely padded copies of $C$.

\subsection{A connected defect in every dimension}

\begin{proof}[Proof of \cref{cor:connected-gap}]
For $n\ge7$, put
\[
v_n=\left\lfloor\frac{n-1}{2}\right\rfloor,
\qquad
u_n=n-v_n-2,
\]
and take the phases
\[
0,1,\ldots,v_n-1.
\]
The inequality $v_n\le u_n+1$ holds, so the snapshot target is defined.  By \cref{thm:snapshot-family},
\[
\ellDown(A_n)=u_n+2v_n=n+v_n-2,
\]
while
\[
\ellAll(A_n)\le u_n+v_n+2=n.
\]
Hence
\[
\ellDown(A_n)-\ellAll(A_n)
\ge v_n-2
=
\left\lfloor\frac{n-5}{2}\right\rfloor.
\]
Connectivity follows from the preceding remark.
\end{proof}

\begin{table}[t]
\centering
\caption{Successive connected cases from the optimized snapshot family.  The restricted value is exact; the ambient column is a certified construction and is exact for $n=7,8,9$ over $\F_2$.}
\label{tab:higher-cases}
\begin{tabular}{@{}rrrrr@{}}
\toprule
$n$&$u$&$v$&$\ellDown$&ambient upper bound / certified gap\\
\midrule
$7$&$2$&$3$&$8$&$7\;/\;1$\\
$8$&$3$&$3$&$9$&$8\;/\;1$\\
$9$&$3$&$4$&$11$&$9\;/\;2$\\
$10$&$4$&$4$&$12$&$10\;/\;2$\\
$11$&$4$&$5$&$14$&$11\;/\;3$\\
$12$&$5$&$5$&$15$&$12\;/\;3$\\
$13$&$5$&$6$&$17$&$13\;/\;4$\\
$15$&$6$&$7$&$20$&$15\;/\;5$\\
$17$&$7$&$8$&$23$&$17\;/\;6$\\
$19$&$8$&$9$&$26$&$19\;/\;7$\\
$21$&$9$&$10$&$29$&$21\;/\;8$\\
\bottomrule
\end{tabular}
\end{table}

\begin{figure}[t]
\centering
\includegraphics[width=.76\linewidth]{connected_higher_cases.pdf}
\caption{The six-row atom is the first exact defect.  From dimension seven onward, the connected snapshot family gives an exact restricted length and a shorter ambient construction in every dimension.}
\label{fig:higher-cases}
\end{figure}

The ratio also increases beyond the block-sum value $7/6$.

\begin{corollary}[Asymptotic multiplicative distortion]\label{cor:three-halves}
For every field,
\[
\limsup_{n\to\infty}
\sup_{A\in\U_n(\F),\,A\ne I}
\frac{\ellDown(A)}{\ellAll(A)}
\ge\frac32.
\]
\end{corollary}

\begin{proof}
Along the optimized connected family,
\[
\frac{\ellDown(A_n)}{\ellAll(A_n)}
\ge
\frac{n+\lfloor(n-1)/2\rfloor-2}{n},
\]
which tends to $3/2$.
\end{proof}

\section{Bounded and weighted catalytic resources}\label{sec:resources}

The connected snapshot family uses a single matched excursion: one upper-directed gate opens the dirty carrier and one upper-directed gate closes it.  This permits a sharper statement than the unrestricted inequality $\ellAll<\ellDown$.

\subsection{A bounded number of upper-directed gates}

For a circuit word $w$, let $u(w)$ denote its number of upper-directed gates.

\begin{definition}[Upper-gate budget]
For $q\ge0$ and $A\in\U_n(\F)$, define
\[
\ellBudget{q}(A)
=
\min\{|w|:[w]=A,\ u(w)\le q\}.
\]
Thus
\[
\ellBudget{0}(A)=\ellDown(A),
\qquad
\ellBudget{q+1}(A)\le\ellBudget{q}(A),
\qquad
\ellAll(A)\le\ellBudget{q}(A).
\]
\end{definition}

\begin{proposition}[Two upper gates give a linear advantage]\label{prop:two-upper-linear}
For every snapshot target $A=A(u;s_1,\ldots,s_v)$,
\[
\ellBudget{0}(A)=u+2v,
\qquad
\ellBudget{2}(A)\le u+v+2.
\]
Consequently,
\[
\ellBudget{0}(A)-\ellBudget{2}(A)\ge v-2.
\]
For the optimized connected target $A_n$ of \cref{cor:connected-gap},
\[
\ellBudget{0}(A_n)-\ellBudget{2}(A_n)
\ge
\left\lfloor\frac{n-5}{2}\right\rfloor.
\]
\end{proposition}

\begin{proof}
The exact restricted value is \cref{thm:snapshot-family}.  Its ambient construction has precisely two upper-directed gates, namely the opening and closing operations on the carrier row; all $u$ updates and all $v$ reads are lower directed.  The optimized choice has $v=\lfloor(n-1)/2\rfloor$.
\end{proof}

\begin{corollary}[No budget-only additive simulation bound]\label{cor:no-budget-only}
There is no finite constant $K$ such that
\[
\ellDown(A)\le\ellBudget{2}(A)+K
\]
for all dimensions, fields, and lower-unitriangular targets $A$.
\end{corollary}

\begin{proof}
The left-minus-right difference in \cref{prop:two-upper-linear} tends to infinity with $n$.
\end{proof}

Thus the power of the excursion is not controlled by the number of forbidden gates.  It is controlled by how many useful carrier states are exported between opening and closing.

\subsection{Rank-one, bandwidth-one flag escape}\label{subsec:microscopic-escape}

The number of upper gates does not measure how far a circuit travels from the lower subgroup.  We now control every intermediate state.

For $X\in\SL_n(\F)$, let $X^+$ be the matrix obtained by retaining only the strict upper entries of $X$.  Define
\[
r^+(X)=\rank X^+
\]
and
\[
b^+(X)=\max\{j-i:X_{ij}\ne0,\ i<j\},
\]
with $b^+(X)=0$ when $X^+=0$.  Let $\ellExc{r}{b}(A)$ be the minimum length of a transvection circuit from $I$ to $A$ such that every prefix matrix $X$ satisfies
\[
r^+(X)\le r,
\qquad
b^+(X)\le b.
\]

\begin{lemma}\label{lem:zero-excursion-equals-lower}
For $A\in\U_n(\F)$,
\[
\ellExc{0}{0}(A)=\ellDown(A).
\]
\end{lemma}

\begin{proof}
A lower-directed gate preserves lower triangularity.  Conversely, if $X$ is lower unitriangular and an upper-directed gate $R_t\leftarrow R_t+aR_s$ with $t<s$ is applied, then the new $(t,s)$ entry equals
\[
X_{t,s}+aX_{s,s}=a\ne0.
\]
The next state is therefore not lower triangular.  A circuit whose every prefix has zero strict upper part uses only lower-directed gates.
\end{proof}

\begin{theorem}[Microscopic escape, macroscopic saving]\label{thm:microscopic-escape}
For every field $\F$ and every $n\ge7$, there is a connected target $A_n\in\U_n(\F)$ such that
\[
\ellExc{0}{0}(A_n)-\ellExc{1}{1}(A_n)
\ge
\left\lfloor\frac{n-5}{2}\right\rfloor.
\]
The length-$n$ circuit witnessing the second term uses exactly two upper-directed gates.
\end{theorem}

\begin{proof}
Use the connected family from \cref{cor:connected-gap}, with the carrier coordinates $p<q$ adjacent.  Between the opening and closing gates, the only strict upper entry is the coefficient at position $(p,q)$.  Hence every prefix has upper rank at most one and upper bandwidth at most one.  The displayed ambient circuit has length $n$, whereas
\[
\ellExc{0}{0}(A_n)=\ellDown(A_n)
=n+\left\lfloor\frac{n-1}{2}\right\rfloor-2.
\]
Subtracting proves the claim.
\end{proof}

\begin{figure}[t]
\centering
\includegraphics[width=.82\linewidth]{microscopic_escape.pdf}
\caption{The snapshot circuit leaves the lower subgroup through one adjacent upper entry.  Throughout the excursion the strict upper part has rank one and bandwidth one, yet the accumulated saving grows linearly.}
\label{fig:microscopic-escape}
\end{figure}

\subsection{A weighted upper-gate metric}

A hardware or fault model may penalize an upper-directed interaction more heavily than a lower-directed one.  Fix $\lambda>0$.  If $w$ has $d(w)$ lower-directed gates and $u(w)$ upper-directed gates, put
\[
\operatorname{cost}_{\lambda}(w)=d(w)+\lambda u(w),
\]
and define
\[
\ellWeight{\lambda}(A)
=
\inf\{\operatorname{cost}_{\lambda}(w):[w]=A\}.
\]
At $\lambda=1$ this is the ordinary ambient length, while large $\lambda$ increasingly penalizes flag excursions.

\begin{proposition}[Fixed penalties do not suppress catalysis]\label{prop:weighted-catalysis}
For every snapshot target $A=A(u;s_1,\ldots,s_v)$ and every $\lambda>0$,
\[
\ellWeight{\lambda}(A)
\le
u+v+2\lambda.
\]
Consequently,
\[
\ellDown(A)-\ellWeight{\lambda}(A)
\ge
v-2\lambda.
\]
For the optimized connected family,
\[
\ellDown(A_n)-\ellWeight{\lambda}(A_n)
\ge
\left\lfloor\frac{n-1}{2}\right\rfloor-2\lambda.
\]
In particular, for every fixed finite $\lambda$, the weighted advantage grows linearly with $n$.
\end{proposition}

\begin{proof}
The displayed catalytic word contains $u+v$ lower-directed gates and two upper-directed gates, so its weighted cost is $u+v+2\lambda$.  Subtract this from the exact restricted value $u+2v$.  The optimized family again has $v=\lfloor(n-1)/2\rfloor$.
\end{proof}

\begin{remark}[Penalty threshold]
For the displayed realization of a snapshot target, catalysis beats the optimal lower-only cost whenever
\[
\lambda<\frac v2.
\]
The critical penalty therefore scales with the number of consumers, not with the dimension of the dirty workspace.
\end{remark}

\begin{figure}[t]
\centering
\includegraphics[width=.78\linewidth]{weighted_catalysis.pdf}
\caption{Certified lower-minus-weighted advantage for the optimized connected family.  Every fixed upper-gate penalty is eventually amortized because the same matched excursion serves a growing number of output snapshots.}
\label{fig:weighted-catalysis}
\end{figure}

\begin{proof}[Proof of \cref{thm:resource-robustness}]
Combine \cref{prop:two-upper-linear,prop:weighted-catalysis}.
\end{proof}


\section{Algebraic rarity of catalytic defects}\label{sec:rarity}

The preceding sections construct highly structured targets with large defects.  This section proves the complementary statement that catalytic defects are nongeneric at every fixed dimension.  Fix a field $K$.  Algebraic dimensions and closures are taken after base change to an algebraic closure of $K$.  Put
\[
N=\dim \U_n=\binom n2,
\qquad
M=n(n-1).
\]

\subsection{Transvection product loci}

For a direction skeleton
\[
\sigma=((t_1,s_1),\ldots,(t_L,s_L)),
\qquad t_i\ne s_i,
\]
define the polynomial word map
\[
\Phi_\sigma:\Aff^L\longrightarrow\SL_n,
\qquad
(a_1,\ldots,a_L)
\longmapsto
(I+a_LE_{t_L,s_L})\cdots(I+a_1E_{t_1,s_1}).
\]
The coefficients are allowed to vanish, so a word of length at most $L$ can be padded to length $L$.  Let
\[
\mathcal R^*_{n,\le L}
=
\U_n\cap
\bigcup_{\sigma\in\{(t,s):t\ne s\}^L}
\operatorname{im}\Phi_\sigma
\]
be the ambient reachability locus, and define $\mathcal R^\downarrow_{n,\le L}$ analogously using only lower directions.

Equivalently, put
\[
X_*
=
\bigcup_{t\ne s}\{I+aE_{t,s}:a\in K\},
\qquad
X_\downarrow
=
\bigcup_{s<t}\{I+aE_{t,s}:a\in K\}.
\]
Then $\ellAll(A)$ and $\ellDown(A)$ are the corresponding transvection product lengths: the least $L$ for which $A\in X_*^L$ or $A\in X_\downarrow^L$.

\begin{proposition}[Dimension of bounded-length loci]\label{prop:reachability-dimension}
For every $L\ge0$, the sets $\mathcal R^*_{n,\le L}$ and $\mathcal R^\downarrow_{n,\le L}$ are constructible and
\[
\dim\overline{\mathcal R^*_{n,\le L}}
\le L,
\qquad
\dim\overline{\mathcal R^\downarrow_{n,\le L}}
\le L.
\]
\end{proposition}

\begin{proof}
There are finitely many direction skeletons of length $L$.  For each skeleton, $\Phi_\sigma$ is a morphism from an $L$-dimensional affine space.  Chevalley's theorem makes its image constructible, and the closure of the image has dimension at most $L$ \cite{Hartshorne1977}.  Finite unions preserve constructibility and the maximum dimension bound.  Intersecting with the closed subvariety $\U_n$ cannot increase dimension.
\end{proof}

The argument is an instance of a general parameter-counting principle.

\begin{proposition}[Dimension per gate]\label{prop:dimension-per-gate}
Let $G$ be an algebraic group, let $V\subseteq G$ be an irreducible subvariety of dimension $D$, and let $X\subseteq G$ be a finite union of locally closed subvarieties, each of dimension at most $r$, with the identity in $X$.  Write
\[
X^{(L)}
=
\{x_L\cdots x_1:x_1,\ldots,x_L\in X\}.
\]
Then
\[
\dim\overline{V\cap X^{(L)}}\le rL.
\]
In particular, if $rL<D$, a Zariski-generic point of $V$ has $X$-product length greater than $L$.
\end{proposition}

\begin{proof}
Every irreducible component of the Cartesian power $X^{\times L}$ has dimension at most $rL$.  Multiplication $G^L\to G$ is a morphism, so the closure of the image of each component has dimension at most $rL$.  There are finitely many components, and intersecting the resulting finite union with $V$ cannot increase dimension.  Since the identity belongs to $X$, products of length at most $L$ are already contained in $X^{(L)}$.
\end{proof}

For transvections, the discrete choice of direction contributes finitely many components and each component carries one scalar parameter.  Thus $r=1$.  Every element of $\U_n(K)$ has lower length at most $N$: process the strict lower entries in any root order compatible with elimination.  The dimension-per-gate bound shows that this worst-case value is generic.

\begin{theorem}[Generic rigidity]\label{thm:generic-rigidity}
Let $K$ be an infinite field.  There is a nonempty Zariski-open subvariety $\Omega_n\subseteq\U_n$, defined over $K$, such that $\Omega_n(K)$ is nonempty and
\[
\ellAll(A)=\ellDown(A)=N
\qquad(A\in\Omega_n(K)).
\]
The same statement holds with exact product length replaced by border product length.
\end{theorem}

\begin{proof}
After base change to an algebraic closure $\overline K$, \cref{prop:reachability-dimension} gives
\[
\dim\overline{\mathcal R^*_{n,\le N-1}}\le N-1<N=\dim\U_n.
\]
Because the word maps are defined over $K$, this proper closed set descends to a proper closed subset of $\U_n$ over $K$.  Its complement is a nonempty Zariski-open set $\Omega_n$.  Since $K$ is infinite and $\U_n\cong\Aff^N$, the set $\Omega_n(K)$ is nonempty and Zariski dense.  For any $A\in\Omega_n(K)$, one has $\ellAll(A)\ge N$.  The inequalities
\[
\ellAll(A)\le\ellDown(A)\le N
\]
force equality.  Replacing each image by its Zariski closure gives the border-length statement by the same argument.
\end{proof}

\subsection{Defect codimension}

For an algebraically closed field $K$ and $d\ge1$, define
\[
\mathcal D_{n,\ge d}
=
\{A\in\U_n(K):\DeltaDef(A)\ge d\}.
\]

\begin{theorem}[Defect--codimension bound]\label{thm:defect-codimension}
Let $K$ be algebraically closed.  For every $1\le d\le N$,
\[
\dim\overline{\mathcal D_{n,\ge d}}\le N-d.
\]
Equivalently,
\[
\codim_{\U_n}\overline{\mathcal D_{n,\ge d}}\ge d.
\]
More finely, on the stratum where $\ellDown(A)=L$ and $\DeltaDef(A)\ge d$, the closure has dimension at most $L-d$.
\end{theorem}

\begin{proof}
If $\DeltaDef(A)\ge d$, then
\[
\ellAll(A)
\le
\ellDown(A)-d
\le
N-d.
\]
Thus $\mathcal D_{n,\ge d}\subseteq\mathcal R^*_{n,\le N-d}$, and \cref{prop:reachability-dimension} gives the first assertion.  On the intrinsic-length-$L$ stratum one instead has $\ellAll(A)\le L-d$, proving the refinement.
\end{proof}

\begin{figure}[t]
\centering
\includegraphics[width=.78\linewidth]{defect_strata.pdf}
\caption{Bounded ambient product loci form nested algebraic strata inside $\U_n$.  A defect of size at least $d$ lies inside the ambient reachability locus of length $N-d$, whose closure has codimension at least $d$.}
\label{fig:defect-strata}
\end{figure}

\subsection{Finite-field sparsity}

The dimension argument has an elementary counting counterpart.  Allowing zero coefficients, a padded word of length $L$ is specified by $L$ choices among $M$ directed positions and $L$ coefficients in $\F_q$.  Consequently
\[
|\mathcal R^*_{n,\le L}(\F_q)|\le(Mq)^L.
\]

\begin{corollary}[Finite-field defect density]\label{cor:finite-field-density}
For $1\le d\le N$,
\[
\Prob_{A\in\U_n(\F_q)}[\DeltaDef(A)\ge d]
\le
\min\{1,M^{N-d}q^{-d}\}.
\]
In particular, for fixed $n$,
\[
\Prob[\DeltaDef(A)>0]=O_n(q^{-1}),
\qquad
\Exp[\DeltaDef(A)]=O_n(q^{-1})
\]
as $q\to\infty$.
\end{corollary}

\begin{proof}
The defect condition implies ambient length at most $N-d$, so the number of such targets is at most $(Mq)^{N-d}$.  Divide by
\[
|\U_n(\F_q)|=q^N.
\]
The expectation bound follows from
\[
\Exp[\DeltaDef]=\sum_{d=1}^{N}\Prob[\DeltaDef\ge d].
\]
\end{proof}

The same counting argument also gives
\[
\Prob_{A\in\U_n(\F_q)}[\ellAll(A)=\ellDown(A)=N]
\ge
1-M^{N-1}q^{-1},
\]
whenever the right side is positive.  Thus maximal ambient and intrinsic length is not only Zariski generic over infinite fields but asymptotically typical over growing finite fields at fixed dimension.

\subsection{Border catalysis}

For an algebraic generating set $X$, define the border product length
\[
\underline\ell_X(A)
=
\min\{L:A\in\overline{X^L}\}.
\]
This is analogous to the passage from tensor rank to border rank \cite{Landsberg2012}.  The generic value is still $N$ by \cref{thm:generic-rigidity}, but exact and border defects can differ on exceptional strata.  It is presently open whether the canonical atom $C$ belongs to the closure of the six-factor lower transvection locus.  A negative answer would show that its one-gate defect is stable under algebraic degeneration; a positive answer would exhibit a rank-versus-border-rank phenomenon inside reversible linear synthesis.

\begin{remark}[Rare but amplifiable]\label{rem:rare-amplifiable}
The codimension theorem and the tensor construction point in opposite directions without contradiction.  At each fixed $n$, a defect of size $d$ lies on a codimension-at-least-$d$ subset.  Carefully selected points on these thin strata can nevertheless be tensorized into families with $\Omega(n\log n)$ additive defect.  Catalytic compression is therefore algebraically rare but asymptotically amplifiable.
\end{remark}

\section{Recursive, block, and tensor amplification}\label{sec:amplification}

\subsection{Restricted block additivity}

The lower metric is exactly additive on contiguous block-diagonal targets.

\begin{proposition}[Block additivity]\label{prop:block-additivity}
Let $A_i\in\U_{n_i}(\F)$ for $1\le i\le k$.  Then
\[
\ellDown\bigl(\operatorname{diag}(A_1,\ldots,A_k)\bigr)
=
\sum_{i=1}^k\ellDown(A_i).
\]
\end{proposition}

\begin{proof}
The upper bound follows by running optimal lower circuits independently inside the blocks.

For the lower bound, view an element of $\U_{n_1+\cdots+n_k}(\F)$ as a lower block-triangular matrix and project to its diagonal blocks:
\[
\pi(X)=(X_{11},\ldots,X_{kk}).
\]
Because diagonal blocks multiply independently in a lower block-triangular product, $\pi$ is a homomorphism.  A lower transvection internal to one block projects to one generator in that factor, while a crossing transvection projects to the identity in every factor.  If $m_i$ gates survive in factor $i$, then the total word length is at least
\[
\sum_i m_i
\ge
\sum_i\ellDown(A_i).
\]
\end{proof}

The ambient metric is subadditive:
\[
\ellAll\bigl(\operatorname{diag}(A_1,\ldots,A_k)\bigr)
\le
\sum_i\ellAll(A_i).
\]
Cross-block ambient workspace can only improve this upper bound.

\begin{corollary}[Block-sum amplification]\label{cor:linear-gap}
For every field and every $n\ge6$, there is an $A_n\in\U_n(\F)$ with
\[
\ellDown(A_n)-\ellAll(A_n)
\ge
\left\lfloor\frac n6\right\rfloor.
\]
For $n=6k$, the target $C^{\oplus k}$ satisfies
\[
\ellDown(C^{\oplus k})=7k,
\qquad
\ellAll(C^{\oplus k})\le6k.
\]
\end{corollary}

\begin{proof}
For $n=6k+r$, take $C^{\oplus k}\oplus I_r$ and apply block additivity together with $k$ copies of the ambient six-gate word.
\end{proof}

\begin{figure}[t]
\centering
\includegraphics[width=.68\linewidth]{direct_sum_amplification.pdf}
\caption{Block-sum amplification.  The restricted length is exactly $7k$, while the displayed ambient construction has length $6k$.}
\label{fig:direct-sum}
\end{figure}

\subsection{A recursive zeta-shaped lift}

For $A\in\U_m(\F)$, define
\begin{equation}\label{eq:zeta-lift}
\mathcal Z(A)
=
\begin{pmatrix}
A&0\\
A&A
\end{pmatrix}
\in\U_{2m}(\F).
\end{equation}

\begin{proposition}[Exact restricted lift]\label{prop:zeta-lift}
For every $A\in\U_m(\F)$,
\[
\ellDown(\mathcal Z(A))=2\ellDown(A)+m,
\]
and
\[
\ellAll(\mathcal Z(A))\le2\ellAll(A)+m.
\]
Consequently,
\[
\DeltaDef(\mathcal Z(A))\ge2\DeltaDef(A).
\]
\end{proposition}

\begin{proof}
The construction first realizes $\operatorname{diag}(A,A)$ and then applies the $m$ root-cross gates that add top row $i$ into bottom row $i$.

For the lower bound, write a circuit state in $m\times m$ blocks as
\[
X=
\begin{pmatrix}
P&0\\
Q&R
\end{pmatrix}.
\]
Projection to the two diagonal blocks forces at least $2\ellDown(A)$ internal gates.  A top-internal gate leaves $Q$ unchanged.  A bottom-internal gate performs an invertible row operation on $Q$ and preserves its rank.  A root-cross gate changes $Q$ by one rank-one row contribution, so it can increase $\rank Q$ by at most one.  The final cross block is $A$, of rank $m$, hence at least $m$ root-cross gates are necessary.
\end{proof}

Let $A^{[0]}=A$ and $A^{[r+1]}=\mathcal Z(A^{[r]})$.  If $A$ has size $m$, restricted length $a$, and ambient length $b$, then
\begin{equation}\label{eq:zeta-iterate}
\ellDown(A^{[r]})
=
2^r a+r m2^{r-1},
\end{equation}
and
\begin{equation}\label{eq:zeta-iterate-ambient}
\ellAll(A^{[r]})
\le
2^r b+r m2^{r-1}.
\end{equation}
Thus
\[
\DeltaDef(A^{[r]})\ge2^r\DeltaDef(A).
\]
For the canonical atom $C$, the first values are
\begin{center}
\begin{tabular}{@{}rrrrr@{}}
\toprule
$r$&dimension&restricted exact&ambient upper bound&certified gap\\
\midrule
$0$&$6$&$7$&$6$&$1$\\
$1$&$12$&$20$&$18$&$2$\\
$2$&$24$&$52$&$48$&$4$\\
$3$&$48$&$128$&$120$&$8$\\
$4$&$96$&$304$&$288$&$16$\\
\bottomrule
\end{tabular}
\end{center}
This gives arbitrarily deep exact recursion with identical diagonal children and a full-rank cross block while preserving a positive defect.

\begin{figure}[t]
\centering
\includegraphics[width=.88\linewidth]{recursive_zeta_lift.pdf}
\caption{The recursive lift has the same outer shape as a Boolean zeta recursion.  A finite defect at the leaf is doubled at every lift.}
\label{fig:zeta-lift}
\end{figure}

\subsection{Tensor products of tree-tight targets}

The tensor operation gives the stronger asymptotic amplification.

\begin{theorem}[Exact restricted tensor rule]\label{thm:tensor-rule}
Let $A\in\U_m(\F)$ and $B\in\U_n(\F)$ be tree tight.  Then $A\otimes B$ is tree tight and
\begin{equation}\label{eq:tensor-lower}
\ellDown(A\otimes B)
=
n\ellDown(A)+m\ellDown(B).
\end{equation}
Moreover,
\begin{equation}\label{eq:tensor-ambient}
\ellAll(A\otimes B)
\le
n\ellAll(A)+m\ellAll(B),
\end{equation}
so
\begin{equation}\label{eq:tensor-defect}
\DeltaDef(A\otimes B)
\ge
n\DeltaDef(A)+m\DeltaDef(B).
\end{equation}
\end{theorem}

\begin{proof}
Use lexicographic order on the $mn$ coordinates.  Substitute a copy of a certificate tree for $B$ at every leaf of a certificate tree for $A$.  At an outer node, the cross block is
\[
A[R,L]\otimes B,
\]
whose rank is $n\rank A[R,L]$ because $B$ is invertible.  The outer nodes therefore contribute $n\ellDown(A)$.  The $m$ diagonal copies of the inner tree contribute $m\ellDown(B)$.  This proves the lower bound.

For the matching lower construction, first realize $I_m\otimes B$ using $m$ copies of a circuit for $B$.  Then realize $A\otimes I_n$.  Each gate of a circuit for $A$ lifts to $n$ coordinatewise gates.  Since
\[
(A\otimes I_n)(I_m\otimes B)=A\otimes B,
\]
the cost is the right side of \cref{eq:tensor-lower}.  The same construction with ambient circuits proves \cref{eq:tensor-ambient}.
\end{proof}

If $A$ is tree tight of dimension $m$, with restricted length $a$ and an ambient construction of length $b$, then
\begin{equation}\label{eq:tensor-power-lower}
\ellDown(A^{\otimes r})
=
r m^{r-1}a,
\end{equation}
and
\begin{equation}\label{eq:tensor-power-upper}
\ellAll(A^{\otimes r})
\le
r m^{r-1}b.
\end{equation}
A positive one-copy defect therefore becomes an order-$N\log N$ defect at dimension $N=m^r$.

\begin{proof}[Proof of \cref{thm:tensor-main}]
Apply \cref{thm:tensor-rule} repeatedly to the tree-tight snapshot target $A=A(u;s_1,\ldots,s_v)$, for which
\[
a=u+2v,
\qquad
b=u+v+2=m.
\]
Then
\[
\ellDown(A^{\otimes r})
-
\ellAll(A^{\otimes r})
\ge
r m^{r-1}(v-2).
\]
Since $N=m^r$ and $r=\log_2N/\log_2m$, this is
\[
\frac{v-2}{m\log_2m}N\log_2N.
\]
\end{proof}

Within the snapshot family, the best coefficient is obtained by the twenty-one-dimensional seed
\[
A_{21}=A(9;0,1,\ldots,9),
\]
for which
\[
\ellDown(A_{21})=29,
\qquad
\ellAll(A_{21})\le21,
\qquad
\DeltaDef(A_{21})\ge8.
\]

\begin{proposition}[Optimal snapshot seed for tensor rate]\label{prop:optimal-snapshot-seed}
Among all snapshot targets, the coefficient
\[
\frac{v-2}{m\log_2m},
\qquad m=u+v+2,
\]
is maximized at
\[
(m,u,v)=(21,9,10).
\]
\end{proposition}

\begin{proof}
For fixed $m$, maximize $v$ subject to $v\le u+1$ and $u+v+2=m$.  Thus
\[
v=\left\lfloor\frac{m-1}{2}\right\rfloor.
\]
If $m$ is even and at least eight, its coefficient is strictly smaller than that of $m-1$, because the numerator is the same while
\[
(m-1)\log(m-1)<m\log m.
\]
It therefore suffices to consider odd $m$.  Apart from the constant change of logarithm base, the odd-dimensional envelope is
\[
f(x)=\frac{x-5}{2x\log x}.
\]
Its derivative has the sign of
\[
5(1+\log x)-x.
\]
For $x>5$ this expression crosses zero exactly once, between $19$ and $21$.  Direct comparison gives $f(21)>f(19)$, after which the envelope decreases.  Hence the unique integer maximizer is $m=21$.
\end{proof}

The first tensor powers are
\begin{center}
\begin{tabular}{@{}rrrrr@{}}
\toprule
$r$&dimension $21^r$&restricted exact&ambient upper bound&certified gap\\
\midrule
$1$&$21$&$29$&$21$&$8$\\
$2$&$441$&$1218$&$882$&$336$\\
$3$&$9261$&$38367$&$27783$&$10584$\\
$4$&$194481$&$1074276$&$777924$&$296352$\\
\bottomrule
\end{tabular}
\end{center}

\begin{figure}[t]
\centering
\includegraphics[width=.74\linewidth]{tensor_amplification.pdf}
\caption{Tensor powers of the twenty-one-dimensional snapshot seed.  Both displayed circuit sizes are of order $N\log N$, and their difference has a positive $N\log N$ coefficient.}
\label{fig:tensor-amplification}
\end{figure}

\subsection{Stable tensor capacity and activation}\label{subsec:stable-tensor}

The tensor inequality also defines a regularized ambient invariant.  Let $A\in\U_m(\F)$ and put
\[
b_r(A)
=
\frac{\ellAll(A^{\otimes r})}{m^{r-1}}.
\]
Applying \cref{eq:tensor-ambient} to $A^{\otimes r}\otimes A^{\otimes s}$ gives
\[
b_{r+s}(A)\le b_r(A)+b_s(A).
\]

\begin{proposition}[Stable tensor length]\label{prop:stable-tensor-length}
The limit
\[
\tau_*(A)
=
\lim_{r\to\infty}
\frac{\ellAll(A^{\otimes r})}{r m^{r-1}}
=
\inf_{r\ge1}
\frac{\ellAll(A^{\otimes r})}{r m^{r-1}}
\]
exists.  If $A$ is tree tight, then the stable catalytic capacity
\[
\kappa_\otimes(A)
=
\ellDown(A)-\tau_*(A)
\]
satisfies
\[
\kappa_\otimes(A)\ge\DeltaDef(A).
\]
\end{proposition}

\begin{proof}
The sequence $b_r(A)$ is subadditive, so Fekete's lemma gives the displayed limit.  Since the one-copy lifted construction implies $\tau_*(A)\le\ellAll(A)$, one has
\[
\kappa_\otimes(A)
\ge
\ellDown(A)-\ellAll(A)
=
\DeltaDef(A).
\]
For a tree-tight target, \cref{eq:tensor-power-lower} identifies $\ellDown(A)$ with its exact normalized restricted tensor cost.
\end{proof}

For the twenty-one-dimensional seed $A_{21}$,
\[
\kappa_\otimes(A_{21})\ge8.
\]
This regularization introduces a new possibility: a target may have $\DeltaDef(A)=0$ but $\kappa_\otimes(A)>0$.  In that case one copy is rigid while many tensor copies admit a collective shortcut.  We call this \emph{tensor activation}.

\subsection{Extremal consequences}

Define
\[
D_n(\F)
=
\sup_{A\in\U_n(\F)}
\bigl(\ellDown(A)-\ellAll(A)\bigr).
\]
For finite fields the supremum is a maximum.  The tensor family gives the following explicit all-dimension bound.

\begin{proof}[Proof of \cref{cor:nlogn-gap}]
Let $m=21$ and write
\[
n=\sum_{j=0}^{r}d_jm^j,
\qquad
0\le d_j<m.
\]
Take the block sum of $d_j$ copies of $A_{21}^{\otimes j}$ for every $j$, with the convention $A_{21}^{\otimes0}=I_1$.  Block additivity and the tensor estimate give
\[
D_n(\F)
\ge
\frac{8}{21}
\sum_{j=0}^{r}j d_j21^j.
\]
Put $S=\sum_j j d_j21^j$.  Since $r=\lfloor\log_{21}n\rfloor$,
\[
r n-S
=
\sum_{j<r}(r-j)d_j21^j
\le
20\sum_{h\ge1}h21^{r-h}
\le
\frac{21^{r+1}}{20}
\le
\frac{21}{20}n.
\]
Also $r\ge\log_{21}n-1$.  Therefore
\[
S\ge n\log_{21}n-\frac{41}{20}n.
\]
Multiplying by $8/21$ yields
\[
D_n(\F)
\ge
\frac{8}{21\log_2 21}n\log_2n
-
\frac{82}{105}n.
\]
\end{proof}

Numerically,
\[
\frac{8}{21\log_2 21}
=0.0867315\ldots.
\]
The elementary bottom-to-top construction gives
\[
D_n(\F)\le\frac{n(n-1)}2.
\]
Thus the current general range is
\begin{equation}\label{eq:defect-range}
0.0867315\,n\log_2n-O(n)
\le
D_n(\F)
\le
\frac{n(n-1)}2.
\end{equation}

\section{Discussion and open problems}\label{sec:discussion}

\subsection{What the threshold and higher families establish}

The six-row atom proves that an optimal circuit for a lower-unitriangular target need not remain in the lower-unitriangular subgroup.  The exact binary classification shows that this first failure is rigid: it occurs for exactly twenty targets, every defect is $(7,6)$, and every target is generated by the same three-cross, one-phase mechanism.

The higher-dimensional results change the scale of the phenomenon.  The connected snapshot family shows that exactly two upper-directed gates can support a defect linear in dimension, and that the advantage persists under every fixed finite penalty on those gates.  The recursive lift shows that a defective leaf survives arbitrarily many identical full-rank dyadic levels.  Tensor products then turn any tree-tight positive-defect atom into an order-$N\log N$ additive separation.  Thus the defect is not a finite anomaly and not merely linear in the worst case.  At the same time, generic rigidity and the defect--codimension theorem show that large defects require correspondingly special algebraic conditions.  The central synthesis is therefore not ubiquity but amplification: thin exceptional strata can control large structured families.

\subsection{Relation to ordinary CNOT synthesis}

Exact CNOT synthesis usually minimizes $\ellAll$ on all of $\GL_n(\F_2)$.  The present problem fixes a structured target class and compares the ambient optimum with a structure-preserving optimum.  Recent exact methods explore the full Cayley graph using symmetries, bidirectional search, A* methods, SAT, and related techniques \cite{ChristensenEtAl2025,WebsterEtAl2025,ShaikVanDePol2025}.  The right-coset computation in this paper is tailored to a subgroup comparison, while the higher-dimensional exact audits use direct bidirectional balls for selected targets.

The 2026 NP-completeness result for all-to-all CNOT distance shows that cancellation and unrestricted parity reuse must be handled in full generality \cite{AcuavivaEtAl2026}.  The present constructions isolate explicit reusable-carrier mechanisms under an ordered restriction.  They do not resolve ambient synthesis complexity.

\subsection{Open directions}

\begin{question}[Universal six-dimensional threshold]
Is six the first defective dimension over every field, or even over every nonzero commutative ring?  The atom gives a universal defective target in dimension six.  Exact no-defect searches below six are proved here for $\F_2$ and $\F_3$.
\end{question}

\begin{question}[Geometry of the defect locus]
For fixed $n$ and $d$, determine the irreducible components and exact dimension of $\overline{\mathcal D_{n,\ge d}}$.  When is the lower bound $\codim\ge d$ sharp?
\end{question}

\begin{question}[Border catalysis]
Does the canonical atom have lower border product length six or seven?  More generally, when can exact catalytic defects disappear under algebraic degeneration?
\end{question}

\begin{question}[Correct growth of the maximum defect]
Is $D_n(\F)$ of order $n\log n$, or can interacting catalysts produce a larger power of $n$?  The current bounds are
\[
\Omega(n\log n)
\le
D_n(\F)
\le
O(n^2).
\]
\end{question}

\begin{question}[Multiplicative distortion]
Is
\[
\sup_{n,A\in\U_n(\F)}
\frac{\ellDown(A)}{\ellAll(A)}
\]
bounded?  The connected snapshot family gives an asymptotic lower bound of $3/2$.  Exceeding $3/2$ appears to require a mechanism beyond one clean carrier exporting independent snapshots.
\end{question}

\begin{question}[Exact ambient complexity of snapshot targets]
For which phase sets is the ambient construction of length $u+v+2$ optimal?  It is exact for the displayed binary targets in dimensions seven, eight, and nine, but the general ambient optimum remains open.
\end{question}

\begin{question}[Tensor synergy and activation]
Can $A\otimes B$ admit an ambient circuit shorter than the lifted bound
\[
n\ellAll(A)+m\ellAll(B)?
\]
Such synergy would only increase the defect, but understanding it is necessary for exact tensor complexity.  Does there exist a tree-tight target with $\DeltaDef(A)=0$ but $\kappa_\otimes(A)>0$?
\end{question}

\begin{question}[Higher catalytic normal forms]
Can interacting or crossing catalytic excursions produce asymptotic behavior impossible for a single clean carrier?  A useful classification would separate laminar, crossing, and singular flag excursions.
\end{question}

\begin{question}[Two-gate multiplicative ceiling]
Among targets realizable with at most two upper-directed gates, is the asymptotic multiplicative distortion at most $3/2$?  The snapshot family reaches this value from below, but the present paper proves no matching upper bound.
\end{question}

\begin{question}[Matroidal classification]
Which regular matroids admit a flagged realization whose shortest ambient state tour is strictly shorter than every flag-compatible lift?  The first atom is the graphic matroid $M(K_4)$ with designated source, carrier, and output roles.
\end{question}

\begin{question}[Weighted phase transition]
For a fixed target family, determine the exact critical penalty $\lambda_c$ at which a flag excursion ceases to improve the optimal weighted cost.  The snapshot construction proves only the lower bound $\lambda_c\ge v/2$ for its displayed carrier word.
\end{question}

\begin{question}[Other root systems]
Let $G$ be a split reductive group with unipotent radical $U$.  Compare word length in positive root subgroups with ambient word length using all root subgroups.  Which root systems admit a catalytic defect, and what is the first rank?
\end{question}

\begin{question}[Certificate extraction]
Can exact ambient lower bounds for large structured targets be accompanied by short human-readable certificates rather than exhaustive balls?  The interval-tree method completely handles the restricted side, while the ambient side still lacks an equally general certificate.
\end{question}


\section{The twenty binary dimension-six defects}\label{app:twenty}

A lower-unitriangular binary matrix is encoded by reading the strict lower entries row by row:
\[
(1,0),
(2,0),(2,1),
\ldots,
(5,0),\ldots,(5,4).
\]
The resulting fifteen-bit integer is denoted by $\operatorname{code}(A)$.  A row is displayed as the integer whose binary expansion is its support vector.  \Cref{tab:twenty} lists all defects.  Every row has $\ellDown=7$ and $\ellAll=6$.

\begin{longtable}{@{}rllllll@{}}
\caption{The twenty defective matrices in $\U_6(\F_2)$.}\label{tab:twenty}\\
\toprule
code&$R_0$&$R_1$&$R_2$&$R_3$&$R_4$&$R_5$\\
\midrule
\endfirsthead
\toprule
code&$R_0$&$R_1$&$R_2$&$R_3$&$R_4$&$R_5$\\
\midrule
\endhead
6585&1&3&4&15&22&38\\
6641&1&3&4&14&23&38\\
6649&1&3&4&15&23&38\\
7601&1&3&4&14&22&39\\
7609&1&3&4&15&22&39\\
7665&1&3&4&14&23&39\\
10014&1&2&7&11&28&41\\
11038&1&2&7&11&28&42\\
12894&1&2&7&11&25&44\\
12958&1&2&7&11&26&44\\
13142&1&2&7&10&29&44\\
13198&1&2&7&9&30&44\\
14102&1&2&7&10&28&45\\
15118&1&2&7&9&28&46\\
17912&1&2&4&15&23&49\\
18936&1&2&4&15&23&50\\
19960&1&2&4&15&23&51\\
20984&1&2&4&15&23&52\\
22008&1&2&4&15&23&53\\
23032&1&2&4&15&23&54\\
\bottomrule
\end{longtable}

For example, the canonical target has row encoding
\[
(1,3,4,14,23,39)
\]
and code $7601$.

\section{Search algorithms and correctness}\label{app:algorithms}

\subsection{Restricted search}

For $\F_2$, each lower-unitriangular matrix is stored as a bit string.  Standard breadth-first search gives exact distance because every edge has unit weight.  The algorithm terminates after visiting all states.

\begin{algorithm}[Restricted breadth-first search]
Initialize $d[I]=0$ and a queue containing $I$.  Repeatedly remove $X$ from the queue.  For every lower transvection $s\to t$, form $Y=\tau_{t,s}(1)X$.  If $Y$ has not been visited, set $d[Y]=d[X]+1$ and append $Y$.  Continue until the queue is empty.
\end{algorithm}

\subsection{Right-coset key}

For an invertible binary matrix $A$, let $c_j$ be column $j$.  Starting at $j=n-1$ and moving left, reduce $c_j$ by the already canonicalized columns $c_{j+1},\ldots,c_{n-1}$, using their pivot coordinates.  The resulting ordered reduced columns form the key $\kappa(A)$.

\begin{lemma}
For $A,B\in\GL_n(\F_2)$,
\[
\kappa(A)=\kappa(B)
\quad\Longleftrightarrow\quad
AU_n=BU_n.
\]
\end{lemma}

\begin{proof}
Right multiplication by $U\in\U_n$ replaces column $j$ by itself plus a combination of columns with index greater than $j$.  Right-to-left reduction is therefore unchanged on a right coset.  Conversely, equal reduced columns imply that the original column frames differ by a unique lower-unitriangular column transformation, so $A=BU$ for some $U\in\U_n$.
\end{proof}

\subsection{Ambient search}

\begin{algorithm}[Radius-four right-coset meet in the middle]
Compute the exact radius-four ball $B_4$ in $\GL_6(\F_2)$.  Bucket its elements by $\kappa$.  For each bucket and every ordered pair $(b,a)$ in that bucket, compute $u=b^{-1}a\in\U_6$ and update
\[
d_*(u)\leftarrow
\min\bigl(d_*(u),d(b)+d(a)\bigr).
\]
\end{algorithm}

Correctness is \cref{lem:mitm-complete}.  The computation also checks that every generated pair lies in $\U_6$ and that every one of the $32768$ targets receives a finite distance.

\subsection{Higher-dimensional bidirectional audits}

For the targets $H_7,H_8,H_9$ of \cref{thm:higher-exact}, the program \texttt{small\_snapshot\_exact.cpp} constructs exact ambient balls from the identity and from the target.  If radius-$r$ and radius-$s$ balls are disjoint, symmetry of the generating set rules out every word of length at most $r+s$.  The tested splits are $3+3$, $3+4$, and $4+4$, respectively.  The source stores matrices in one 128-bit word and uses exact row-addition transitions; no probabilistic or floating-point step occurs.

\subsection{Symbolic higher-family checks}

The program \texttt{higher\_family\_tests.py} generates the snapshot targets over $\F_2$, $\F_3$, and $\F_5$, evaluates both displayed circuits, computes the interval-tree potentials, audits a tensor certificate for the seven-dimensional seed, verifies the recursive cross rank, and reproduces the twenty-one-dimensional rate optimization.

\subsection{Further generalization audits}

The program \texttt{further\_generalization\_tests.py} verifies the universal circuits over $\mathbb Z/4\mathbb Z$, $\mathbb Z/6\mathbb Z$, and $\mathbb Z/8\mathbb Z$, audits the rank-one bandwidth-one excursion, checks the finite-field counting bounds symbolically, and verifies the stable tensor subadditivity identities.

\subsection{Motif comparison}

The motif program independently loops over the three adjacent catalyst positions, the phase orientation, the phase leaf, and every nonempty proper before/after subset.  It evaluates each six-gate word from the identity and stores its target.  The resulting set has size twenty.  Exact comparison with the twenty census defects returns equality.

\section{Reproducibility ledger}\label{app:reproducibility}

The ancillary directory \texttt{supplement/} contains portable source, frozen expected outputs, and a one-command driver:
\begin{verbatim}
cd supplement
./run_all.sh
\end{verbatim}
The driver compiles every C++ search, executes the Python motif, higher-family, matroid, and weighted-resource checks, and compares all generated outputs byte for byte with the frozen records.  The supplement README gives the individual commands and the role of every program.  A SHA-256 manifest covers every included source and expected-output file.

\section*{Acknowledgement of computational status}

All theorem statements are labeled by proof type in the body.  The universal ring-valued atom, Schreier-holonomy interpretation, regular-matroid identification, interval-tree family, bounded, weighted, and microscopic resource theorems, algebraic-rarity bounds, block additivity, recursive lift, tensor rule, stable tensor limit, and asymptotic amplification are proved symbolically.  The first-threshold classification, finite censuses, and exact ambient values in dimensions seven through nine are computer assisted, with complete search-space coverage justified in \cref{app:algorithms}.

\section*{Data and code availability}

The source archive accompanying this paper contains the complete LaTeX source, all vector figures, every C++ and Python verification program, frozen expected outputs, and SHA-256 manifests.  The command
\begin{verbatim}
cd supplement
./run_all.sh
\end{verbatim}
rebuilds the finite searches and symbolic audits and compares every output byte for byte with its frozen record.  No randomized or floating-point decision enters a theorem certificate.

